\subsection{Factorialité des anneaux locaux réguliers}
\label{subsection:IV.21.11}
\label{IV.21.11}

\begin{theorem}[21.11.1]
\label{IV.21.11.1}
\textup{(Auslander-Buchsbaum)} ---
Un anneau local noethérien régulier est factoriel.
\end{theorem}

\begin{proof}
La démonstration qui suit est due à I.~Kaplansky.

Soit $A$ un anneau local noethérien régulier de dimension $n$; nous allons raisonner
par récurrence sur $n$. Pour $n=0$, $A$ est un corps et pour $n=1$, un anneau de
valuation discrète, donc principal (et \emph{a fortiori} factoriel). Supposons donc
$n\geq2$ et le théorème démontré pour des anneaux réguliers de dimension $<n$.
Posons $X=\operatorname{Spec}(A)$, désignons par $a$ le point fermé de $X$ et posons
$U=X-\{a\}$. En tout point $y\in U$, on a
$\dim(\mathcal O_{X,y})\leq n-1$, et puisque $A$ est régulier, les anneaux
$\mathcal O_{X,y}$ le sont aussi $(\mathbf 0, 17.3.2)$, donc l'hypothèse de
récurrence entraîne qu'ils sont factoriels. De plus on a
$\operatorname{prof}(A)=\dim(A)\geq2$ puisque $A$ est régulier, donc de
Cohen--Macaulay $(\mathbf 0, 17.1.3)$.
\oldpage[IV-4]{303}
Utilisant (21.6.14), on est ramené à prouver que $\operatorname{Pic}(U)=0$.
Considérons donc un $\mathcal O_U$-Module inversible $\mathscr L$, et prouvons qu'il
est isomorphe à $\mathcal O_U$. Il résulte de $(\mathbf I, 9.4.5)$ qu'il existe un
$\mathcal O_X$-Module \emph{cohérent} $\mathscr F$ tel que
$\mathscr F|U=\mathscr L$. Comme $A$ est régulier, donc de dimension cohomologique
finie $(\mathbf 0, 17.3.1)$, il existe une résolution gauche \emph{finie} de
$\mathscr F$ :
\[
  0\leftarrow\mathscr F\leftarrow\mathcal O_X^{n_1}
  \leftarrow\mathcal O_X^{n_2}\leftarrow\cdots
  \leftarrow\mathcal O_X^{n_h}\leftarrow0
\]
$(\mathbf 0, 17.2.8$ et $\mathbf 0, 17.2.2$, (iii)). Par restriction à $U$, on a
donc une résolution finie
\[
\label{IV.21.11.1.1}
  0\leftarrow\mathscr L\leftarrow\mathcal O_U^{n_1}
  \leftarrow\mathcal O_U^{n_2}\leftarrow\cdots
  \leftarrow\mathcal O_U^{n_h}\leftarrow0.
  \tag{21.11.1.1}
\]
\end{proof}

Le théorème résultera alors des considérations générales suivantes. Sur un espace
annelé $X$, soit $\mathscr E$ un $\mathcal O_X$-Module localement libre de rang fini;
on désignera par $\Lambda^{\max}\mathscr E$ le $\mathcal O_X$-Module inversible qui,
au voisinage de chaque point de $X$, est égal à $\Lambda^p\mathscr E$, en désignant
par $p$ le rang de $\mathscr E$ dans ce voisinage (qui peut varier avec la composante
connexe de $X$). Avec cette notation, on a le

\begin{lemma}[21.11.1.2]
\label{IV.21.11.1.2}
Soient $X$ un espace annelé en anneaux locaux, et
\[
  0\leftarrow\mathscr E_0\xleftarrow{u_0}\mathscr E_1
  \leftarrow\cdots\leftarrow\mathscr E_h\leftarrow0
\]
une suite exacte de $\mathcal O_X$-Modules localement libres de rang fini; alors le
$\mathcal O_X$-Module inversible
\[
  \bigotimes_{0\leq i\leq h}
  (\Lambda^{\max}\mathscr E_i)^{\otimes(-1)^i}
\]
est isomorphe à $\mathcal O_X$.
\end{lemma}

\begin{proof}
Montrons comment ce lemme permettra de conclure la preuve de (21.11.1). Il suffit de
noter pour cela que, pour tout entier $n$,
$\Lambda^{\max}(\mathcal O_U^n)=\Lambda^n\mathcal O_U^n$ est isomorphe à
$\mathcal O_U$, ainsi que $\mathcal O_U^{\otimes(-1)}$. Comme d'autre part
$\Lambda^{\max}\mathscr L=\mathscr L$ pour tout $\mathcal O_U$-Module inversible
$\mathscr L$, le lemme (21.11.1.2), appliqué à la suite exacte (21.11.1.1) montre
que $\mathscr L$ est isomorphe à $\mathcal O_U$.

Reste à prouver (21.11.1.2); on procède par récurrence sur $h$, le lemme étant trivial
pour $h=1$. Pour $h>1$, $\mathscr N=\operatorname{Ker}(u_0)$ est un
$\mathcal O_X$-Module localement libre de rang fini $(\mathbf 0_{\mathrm I}, 5.5.5)$,
et on a les deux suites exactes
\[
  0\leftarrow\mathscr E_0\leftarrow\mathscr E_1
  \leftarrow\mathscr N\leftarrow0
\]
\[
  0\leftarrow\mathscr N\leftarrow\mathscr E_2
  \leftarrow\cdots\leftarrow\mathscr E_h\leftarrow0.
\]
En vertu de l'hypothèse de récurrence,
$(\Lambda^{\max}\mathscr N)\otimes
(\bigotimes_{2\leq i\leq h}(\Lambda^{\max}\mathscr E_i)^{\otimes(-1)^{i-1}})$
est isomorphe à $\mathcal O_X$; il suffira donc de définir un isomorphisme canonique
\[
  (\Lambda^{\max}\mathscr N)\otimes(\Lambda^{\max}\mathscr E_0)
  \xrightarrow{\sim}\Lambda^{\max}\mathscr E_1
\]
pour achever la démonstration. Or, il existe un recouvrement ouvert $(U_\alpha)$ de
$X$ tel que dans tout $U_\alpha$, $\mathscr E_1|U_\alpha$ soit somme directe de
$\mathscr N|U_\alpha$ et d'un $\mathcal O_{U_\alpha}$-Module localement libre
$\mathscr M_\alpha$ $(\mathbf 0_{\mathrm I}, 5.5.5)$, d'où un isomorphisme canonique
$v_\alpha:\mathscr M_\alpha\xrightarrow{\sim}\mathscr E_0|U_\alpha$.
Comme on a un isomorphisme canonique
\[
  r_\alpha:(\Lambda^{\max}\mathscr N|U_\alpha)
  \otimes(\Lambda^{\max}\mathscr M_\alpha)
  \xrightarrow{\sim}(\Lambda^{\max}\mathscr E_1)|U_\alpha,
\]
on en déduit au moyen de $v_\alpha$ un isomorphisme
\[
  u_\alpha:(\Lambda^{\max}\mathscr N|U_\alpha)
  \otimes(\Lambda^{\max}\mathscr E_0|U_\alpha)
  \xrightarrow{\sim}\Lambda^{\max}\mathscr E_1|U_\alpha.
\]
\oldpage[IV-4]{304}
Et il s'agit de montrer que $u_\alpha$ et $u_\beta$ coïncident dans
$U_\alpha\cap U_\beta$ pour deux indices quelconques $\alpha$, $\beta$. Or, si
$v'_\alpha$ et $v'_\beta$ sont les restrictions à $U_\alpha\cap U_\beta$ de
$v_\alpha$ et $v_\beta$ respectivement, on a
$v'_\alpha=v'_\beta\circ w_{\beta\alpha}$, où
\[
  w_{\beta\alpha}:\mathscr M_\alpha|(U_\alpha\cap U_\beta)
  \xrightarrow{\sim}\mathscr M_\beta|(U_\alpha\cap U_\beta)
\]
est la « projection parallèle à $\mathscr N$ » telle que pour toute section
$s\in\Gamma(U_\alpha\cap U_\beta,\mathscr M_\alpha)$,
$w_{\beta\alpha}(s)=s+t$ avec
$t\in\Gamma(U_\alpha\cap U_\beta,\mathscr N)$; l'identité de $u_\alpha$ et de
$u_\beta$ résulte aussitôt de ce fait et de la définition de l'isomorphisme canonique
$r_\alpha$ (Bourbaki, \emph{Alg.}, chap.~III, 3\textsuperscript{e} éd.).
\end{proof}

\subsection{Le théorème de pureté de Van der Waerden pour l'ensemble de
ramification d'un morphisme birationnel}
\label{subsection:IV.21.12}
\label{IV.21.12}

\begin{env}[21.12.1]
\label{IV.21.12.1}
Soient $X$ et $U$ deux préschémas, $f:U\to X$ un morphisme quasi-compact et
quasi-séparé, de sorte que $f_*(\mathcal O_U)$ est une $\mathcal O_X$-Algèbre
quasi-cohérente (1.7.4). On appelle \emph{enveloppe affine} du $X$-préschéma $U$ le
$X$-préschéma affine sur $X$
\[
  U^0=\operatorname{Aff}(U/X)=\operatorname{Spec}(f_*(\mathcal O_U))
  =\operatorname{Spec}(\mathscr A(U))\qquad(\mathbf{II}, 1.1.1).
\]
Si $f^0:U^0\to X$ est le morphisme structural, on a donc par définition
\[
  \mathscr A(U^0)=f^0_*(\mathcal O_{U^0})=f_*(\mathcal O_U)=\mathscr A(U),
\]
et à l'isomorphisme identique de $f_*(\mathcal O_U)$ il correspond par
$(\mathbf{II}, 1.2.7)$ un $X$-morphisme canonique
\[
\label{IV.21.12.1.1}
  i_U:U\to U^0.
  \tag{21.12.1.1}
\]

Pour que $i_U$ soit un \emph{isomorphisme}, il faut et il suffit évidemment que le
morphisme $f:U\to X$ soit \emph{affine}. Pour tout $X$-préschéma $V$ \emph{affine sur
$X$}, l'application $u\mapsto u\circ i_U$ est une \emph{bijection}
\[
  \operatorname{Hom}_X(U^0,V)\xrightarrow{\sim}\operatorname{Hom}_X(U,V)
\]
fonctorielle en $V$: cela résulte de l'existence des bijections canoniques
\[
  \operatorname{Hom}_X(U,V)\xrightarrow{\sim}
  \operatorname{Hom}_{\mathcal O_X}(\mathscr A(V),\mathscr A(U))
\]
et
$\operatorname{Hom}_X(U^0,V)\xrightarrow{\sim}
\operatorname{Hom}_{\mathcal O_X}(\mathscr A(V),\mathscr A(U^0))$
$(\mathbf{II}, 1.2.7)$.

On peut donc dire que $U^0$ \emph{représente} le foncteur covariant
$V\mapsto\operatorname{Hom}_X(U,V)$ dans la catégorie des $X$-préschémas affines sur
$X$ $(\mathbf 0_{\mathrm{III}}, 8.1.11)$. On en déduit
$(\mathbf 0_{\mathrm{III}}, 8.1.7)$ que pour $X$ fixé,
$U\mapsto\operatorname{Aff}(U/X)$ est un foncteur covariant de la catégorie des
$X$-préschémas quasi-compacts et quasi-séparés sur $X$, dans la catégorie des
$X$-préschémas affines sur $X$; de façon plus précise, si $U_1$, $U_2$ sont deux
$X$-préschémas quasi-compacts et quasi-séparés sur $X$, à tout $X$-morphisme
$g:U_1\to U_2$ correspond l'unique $X$-morphisme $g^0:U_1^0\to U_2^0$ rendant
commutatif le diagramme
\[
  \xymatrix{
    U_1 \ar[r]^g \ar[d]_{i_{U_1}} & U_2 \ar[d]^{i_{U_2}} \\
    U_1^0 \ar[r]_{g^0} & U_2^0.
  }
\]
\end{env}

\oldpage[IV-4]{305}
Plus généralement, considérons un diagramme commutatif
\[
  \xymatrix{
    U \ar[r]^v \ar[d]_f & U' \ar[d]^{f'} \\
    X \ar[r]_u & X'
  }
\]
où les morphismes $f$, $f'$ sont quasi-compacts et quasi-séparés et le morphisme $u$
\emph{affine}. Si l'on pose $h=u\circ f$, on a
$h_*(\mathcal O_U)=u_*(f_*(\mathcal O_U))=u_*(f^0_*(\mathcal O_{U^0}))$ et
$u\circ f^0$ est un morphisme affine; on a par suite
$U^0=\operatorname{Aff}(U/X')$ (relativement au morphisme $h$), d'où un unique
$X'$-morphisme $v^0:U^0\to U'^0$ rendant commutatif le diagramme
\[
  \xymatrix{
    U \ar[r]^v \ar[d]_{i_U} & U' \ar[d]^{i_{U'}} \\
    U^0 \ar[r]_{v^0} & U'^0.
  }
\]

\begin{proposition}[21.12.2]
\label{IV.21.12.2}
Soient $f:U\to X$ un morphisme quasi-compact et quasi-séparé,
$h:X'\to X$ un morphisme plat; on pose $U'=U\times_X X'$,
$f'=f_{\{X'\}}:U'\to X'$. Alors on a un $X'$-isomorphisme canonique
\[
\label{IV.21.12.2.1}
  \operatorname{Aff}(U'/X')\xrightarrow{\sim}
  \operatorname{Aff}(U/X)\times_X X'.
  \tag{21.12.2.1}
\]
\end{proposition}

\begin{proof}
En effet, on a
$\operatorname{Aff}(U'/X')=\operatorname{Spec}(f'_*(\mathcal O_{U'}))$, et
$\operatorname{Aff}(U/X)\times_X X'=\operatorname{Spec}(h^*(f_*(\mathcal O_U)))$
$(\mathbf{II}, 1.5.1)$; l'isomorphisme de l'énoncé provient de l'isomorphisme canonique
$h^*(f_*(\mathcal O_U))\xrightarrow{\sim}f'_*(\mathcal O_{U'})$ (2.3.1).
\end{proof}

\begin{corollary}[21.12.3]
\label{IV.21.12.3}
Pour tout morphisme quasi-compact et quasi-séparé $f:U\to X$ et tout $x\in X$, on a,
à un isomorphisme canonique près
\[
  U^0\times_X\operatorname{Spec}(\mathcal O_{X,x})
  =(U\times_X\operatorname{Spec}(\mathcal O_{X,x}))^0.
\]
\end{corollary}

Il résulte aussi de (21.12.2) que l'on a, à un isomorphisme canonique près, pour tout
ouvert $V$ de $X$
\[
\label{IV.21.12.4}
  (f^{-1}(V))^0=(f^0)^{-1}(V).
  \tag{21.12.4}
\]

\begin{env}[21.12.5]
\label{IV.21.12.5}
Nous allons considérer en particulier le cas où $f:U\to X$ est une \emph{immersion
ouverte}, $U$ s'identifiant donc à un préschéma induit sur un ouvert de $X$. Comme le
morphisme $f^{-1}(U)\to U$ est l'identité, il résulte de (21.12.4) que le morphisme
$(f^0)^{-1}(U)\to U$ restriction de $f^0$ est un isomorphisme, $i_U:U\to U^0$ étant
donc une \emph{immersion ouverte} permettant d'identifier $U$ à un préschéma induit
sur un ouvert de $U^0$.

Plus généralement, pour un ouvert $V$ de $X$, la restriction $(f^0)^{-1}(V)\to V$ de
$f^0$ est un isomorphisme de $(f^0)^{-1}(V)$ sur le préschéma induit sur l'ouvert
$V\cap U$ si et seulement si l'immersion ouverte $U\cap V\to V$ est un morphisme
\emph{affine}. Il est clair $(\mathbf{II}, 1.2.1)$
\oldpage[IV-4]{306}
que la réunion de ces ouverts $V$ est le \emph{plus grand} d'entre eux, $U_1$, qui
contient $U$; en vertu de ce qui précède, $U_1$ est aussi le plus grand ouvert ne
rencontrant pas l'ensemble
\[
  \operatorname{Daf}(U/X)=f^0(U^0)-U
\]
(« défaut d'affinité » de l'ouvert $U$ relativement à $X$, qui est vide si et seulement
si $U$ est affine sur $X$); autrement dit, l'ensemble fermé $Z=X-U_1$ est
l'\emph{adhérence} de l'ensemble $\operatorname{Daf}(U/X)$.

On notera que pour tout morphisme \emph{plat} $h:X'\to X$, si l'on pose
$U'=h^{-1}(U)$, on a
\[
\label{IV.21.12.5.1}
  \operatorname{Daf}(U'/X')=h^{-1}(\operatorname{Daf}(U/X))
  \tag{21.12.5.1}
\]
comme il résulte aussitôt de (21.12.2.1) et de $(\mathbf I, 3.4.8)$. En particulier,
pour tout ouvert $V$ de $X$, on a
\[
\label{IV.21.12.5.2}
  \operatorname{Daf}((U\cap V)/V)=\operatorname{Daf}(U/X)\cap V
  \tag{21.12.5.2}
\]
et pour tout $x\in X$,
\[
\label{IV.21.12.5.3}
  \operatorname{Daf}((U\cap\operatorname{Spec}(\mathcal O_{X,x}))
  /\operatorname{Spec}(\mathcal O_{X,x}))
  =\operatorname{Daf}(U/X)\cap\operatorname{Spec}(\mathcal O_{X,x}).
  \tag{21.12.5.3}
\]
\end{env}

Nous allons, lorsque $X$ est localement noethérien, donner des précisions sur la nature
de l'ensemble $\operatorname{Daf}(U/X)$, qui montreront par exemple que lorsque $U$ est
partout dense dans $X$, $\operatorname{Daf}(U/X)$ n'est pas un ensemble fermé rare
quelconque :

\begin{theorem}[21.12.6]
\label{IV.21.12.6}
Soient $X$ un préschéma localement noethérien, $U$ un ouvert non vide de $X$,
$f:U\to X$ l'injection canonique. Alors :
\begin{enumerate}[label=(\roman*)]
  \item L'adhérence $Z=X-U_1$ de $\operatorname{Daf}(U/X)$ est de codimension
  $\geq2$ dans $X$.
  \item Si $T=X-U\supset Z$ est de codimension $\geq2$, le morphisme
  $f^0:U^0\to X$ est surjectif.
\end{enumerate}
\end{theorem}

\begin{proof}
(i) Montrons d'abord que pour tout point $x\in\operatorname{Daf}(U/X)$ on a
nécessairement $\dim(\mathcal O_{X,x})>1$. En effet, $x$ ne peut évidemment être point
maximal de $X$, étant contenu dans $\overline U-U$; il s'agit donc de voir que l'on ne
peut avoir $\dim(\mathcal O_{X,x})=1$. Or, cette relation entraînerait, par (21.12.5.3),
$x\in\operatorname{Daf}(U^{(x)}/\operatorname{Spec}(\mathcal O_{X,x}))$, où
$U^{(x)}=U\cap\operatorname{Spec}(\mathcal O_{X,x})$. Mais les seuls ouverts de
$\operatorname{Spec}(\mathcal O_{X,x})$ sont $\operatorname{Spec}(\mathcal O_{X,x})$
lui-même et les parties de l'ensemble (fini) des points maximaux de
$\operatorname{Spec}(\mathcal O_{X,x})$. Or, l'ensemble ouvert réduit à un point maximal
de $\operatorname{Spec}(\mathcal O_{X,x})$ est affine, par définition des préschémas; on
en conclut que \emph{tous} les ensembles ouverts dans
$\operatorname{Spec}(\mathcal O_{X,x})$ sont affines, donc
$\operatorname{Daf}(U^{(x)}/\operatorname{Spec}(\mathcal O_{X,x}))=\varnothing$
contrairement à l'hypothèse faite.

Pour prouver (i) il faut montrer davantage, savoir que si $x\in X$ est tel que
$\dim(\mathcal O_{X,x})=1$, il existe un voisinage ouvert $W$ de $x$ dans $X$ tel que
$W$ ne rencontre pas $\operatorname{Daf}(U/X)$, c'est-à-dire tel que l'injection
canonique $f_W:U\cap W\to W$ soit affine. Mais, avec les mêmes notations que ci-dessus,
on vient de voir que l'injection canonique
$f^{(x)}:U^{(x)}\to\operatorname{Spec}(\mathcal O_{X,x})$ est affine. On peut
évidemment se borner au cas où $X$ est noethérien; comme le morphisme $f$ est de
présentation finie, la conclusion résulte de (8.10.5, (viii)) appliqué suivant la
méthode de (8.1.2, a)).

\oldpage[IV-4]{307}
(ii) Il s'agit de prouver que pour tout point $x\in T$, on a $x\in f^0(U^0)$; nous
allons montrer d'abord qu'on peut se réduire au cas où $X=\operatorname{Spec}(A)$, où
$A$ est un anneau local noethérien \emph{complet}, et $x$ le point fermé de $X$. Pour
cela, il suffit de faire le changement de base
$h:X'=\operatorname{Spec}(\widehat{\mathcal O}_{X,x})\to X$; si l'on pose
$U'=h^{-1}(U)$, $f'=f_{\{X'\}}$ est l'injection canonique $U'\to X'$, et puisque le
morphisme $h$ est plat, il résulte de (21.12.2) que si l'on prouve que $x$ appartient à
$f'^{,0}(U'^{,0})$, on en déduira que $x\in f^0(U^0)$. En vertu de (6.1.1), on a
bien opéré la réduction cherchée.

Soit alors $X_1$ un sous-préschéma fermé réduit de $X$ ayant pour espace sous-jacent
une composante irréductible de $X$, de dimension \emph{maximale}, parmi celles qui
contiennent une composante irréductible de $T$, et posons
$U_1=U\cap X_1$, $T_1=T\cap X_1=X_1-U_1$; on a donc
$\operatorname{codim}(T_1,X_1)\geq2$ $(\mathbf 0, 14.2.1)$, et le couple $(U_1,X_1)$
vérifie donc les mêmes hypothèses que le couple $(U,X)$ ($X_1$ étant spectre d'un
anneau quotient de $A$, donc local noethérien et complet). On a par ailleurs un
diagramme commutatif
\[
  \xymatrix{
    U_1 \ar[r]^j \ar[d]_{f_1} & U \ar[d]^f \\
    X_1 \ar[r]_i & X
  }
\]
où $i$ et $j$ sont les injections canoniques, $i$ étant donc un morphisme affine; on en
déduit (21.12.1) l'existence d'un morphisme $j^0:U_1^0\to U^0$ rendant commutatif le
diagramme
\[
  \xymatrix{
    U_1^0 \ar[r]^{j^0} \ar[d]_{f_1^0} & U^0 \ar[d]^{f^0} \\
    X_1 \ar[r]_i & X
  }
\]
et par suite, pour prouver que $x\in f^0(U^0)$, il suffit de prouver que
$x\in f_1^0(U_1^0)$. On peut donc remplacer $X$ par $X_1$, et on peut par suite
supposer que l'anneau $A$ est en outre \emph{intègre}. Mais $A$ est quotient d'un
anneau noethérien régulier en vertu du théorème de Cohen $(\mathbf 0, 19.8.8$, (i)) et
puisqu'il est intègre, on peut appliquer (5.11.1) avec la famille $(x_\alpha)$ réduite au
seul point maximal de $X$; comme $\operatorname{codim}(T,X)\geq2$ par hypothèse, on
voit que $f_*(\mathcal O_U)$ est un $\mathcal O_X$-Module \emph{cohérent}, donc le
morphisme $f^0:U^0\to X$ est fini; comme ce morphisme est dominant ($U$ étant partout
dense), il est surjectif $(\mathbf{II}, 6.1.10)$, et par suite $x\in f^0(U^0)$. C.Q.F.D.
\end{proof}

\begin{corollary}[21.12.7]
\label{IV.21.12.7}
Soient $X$ un préschéma localement noethérien, $U$ un sous-préschéma induit sur un
ouvert de $X$, $j:U\to X$ l'immersion canonique; supposons que $j$ soit un morphisme
affine. Alors toute composante irréductible de $T=X-U$ est de codimension $\leq1$
(et par suite de codimension $1$ si $U$ est partout dense).
\end{corollary}

\begin{proof}
Supposons en effet qu'une des composantes irréductibles $T_1$ de $T$ soit de
codimension $\geq2$ dans $X$. Remplaçant au besoin $X$ par un voisinage ouvert du
point générique de $T_1$, on peut supposer $T$ irréductible et de codimension $\geq2$.
Mais alors l'hypothèse
\oldpage[IV-4]{308}
que $j:U\to X$ est affine implique que $U^0$ s'identifie à $U$ et $j^0$ à $j$,
autrement dit que $j^0(U^0)=U$; or, cela contredit la conclusion de (21.12.6) qui,
sous l'hypothèse de dimension faite implique que $j^0$ doit être surjectif.
\end{proof}

\begin{env}[21.12.8]
\label{IV.21.12.8}
Soient $A$ un anneau local noethérien, $Y=\operatorname{Spec}(A)$, $y$ l'unique point
fermé de $Y$, $Y'=Y-\{y\}$. Considérons la condition suivante :
\begin{itemize}
  \item[\textup{(W)}] \emph{Pour tout ouvert $U$ de $Y$ contenu dans $Y'$, ne
  contenant aucune composante irréductible de $Y'$ et tel que l'immersion canonique
  $U\to Y'$ soit affine, $U$ est lui-même un ouvert affine.}
\end{itemize}
Cette condition est entraînée par la suivante :
\begin{itemize}
  \item[\textup{(W')}] \emph{Pour toute partie fermée $T$ de $Y$ dont toute composante
  irréductible est de codimension $1$ dans $Y$, et tel que pour toute composante
  irréductible $Y_i$ de $Y$, on ait $Y_i\cap T\ne\{y\}$, l'ouvert $U=Y-T$ est affine.}
\end{itemize}
En effet, si (W') est vérifiée et si l'ouvert $U$ de $Y$ vérifie l'hypothèse de (W), il
résulte de (21.12.7) qu'aucune composante irréductible de $T=Y-U$ ne peut être de
codimension $\geq2$; comme par ailleurs $U$ ne contient aucune des composantes
irréductibles $Y_i-\{y\}$ de $Y'$, la condition (W') montre que $U$ est affine.

On notera que la condition (W') se simplifie lorsque $Y$ est \emph{irréductible} et est
alors équivalente à la suivante :
\begin{itemize}
  \item[\textup{(W'')}] \emph{Pour toute partie fermée irréductible $T$ de $Y$, de
  codimension $1$ dans $Y$, $Y-T$ est un ouvert affine.}
\end{itemize}
En effet, il est clair que (W') entraîne (W'') lorsque $Y$ est irréductible, et la
réciproque est vraie en considérant les composantes irréductibles de $T$ et en notant
que l'intersection d'un nombre fini d'ouverts affines est un ouvert affine
$(\mathbf I, 5.5.6)$.
\end{env}

\begin{env}[21.12.9]
\label{IV.21.12.9}
\emph{Exemples}. ---
Si $A$ est un anneau local noethérien \emph{factoriel}, il vérifie (W') (et
\emph{a fortiori} (W)), puisque tout idéal premier de hauteur $1$ est principal
$(\mathbf I, 1.3.6)$. Mais il y a des anneaux locaux noethériens non factoriels
vérifiant (W''), par exemple ceux de dimension $\leq1$: en effet, on a remarqué dans la
démonstration de (21.12.6, (i)) que \emph{tous} les ouverts de $Y$ sont alors affines.
On peut d'autre part prouver en utilisant la théorie de la dualité locale (chap.~III,
3\textsuperscript{e} partie) que tout anneau local noethérien de dimension $2$ vérifie
(W').
\end{env}

L'intérêt de la condition (W) réside dans le résultat suivant :

\begin{proposition}[21.12.10]
\label{IV.21.12.10}
Soient $X$, $Y$ deux préschémas localement noethériens, $g:X\to Y$ un morphisme,
$y$ un point fermé de $Y$, $Y'=Y-\{y\}$, $X'=g^{-1}(Y')$; supposons que le morphisme
$g':X'\to Y'$ restriction de $g$ soit une immersion ouverte, et que l'anneau local
$\mathcal O_{Y,y}$ vérifie la condition (W) (21.12.8). Alors, pour toute composante
irréductible $Z$ de $X_y=g^{-1}(y)$, ou bien $Z$ est de codimension $\leq1$ dans $X$,
ou bien son point générique est isolé dans $Z$. Si $Z$ est localement de type fini sur
$k(y)$, la seconde alternative implique que $Z$ est réduit à un seul point.
\end{proposition}

\begin{proof}
La dernière assertion de l'énoncé résulte de ce que $Z$ est alors un préschéma de
Jacobson (10.4.7), et comme l'ensemble des points fermés de $Z$ est alors dense dans
$Z$, le point générique de $Z$ ne peut être isolé dans $Z$ que si $Z$ est réduit à un
seul point.

\oldpage[IV-4]{309}
Supposons qu'il existe dans $X_y$ une composante irréductible $Z$ dont le point
générique $z$ n'est pas isolé dans $Z$, et telle que $\operatorname{codim}(Z,X)\geq2$.
La question étant locale sur $X$ et $Y$ puisque $z$ est non isolé dans $Z$, on peut, en
remplaçant $X$ par un voisinage ouvert de $z$ dans $X$, supposer $X$ et $Y$ affines,
$X_y=Z$ irréductible, non réduit à un point et telle que
$\operatorname{codim}(X_y,X)\geq2$. L'image $g(X-X_y)=U$ est un ouvert de $Y'$
isomorphe à $X-X_y$ par hypothèse; montrons qu'en remplaçant au besoin $X$ par un
voisinage ouvert de $z$ dans $X$, on peut supposer que $U$ ne contient aucune des
composantes irréductibles de $Y'$ dont l'adhérence contient $y$.

En effet, on peut tout d'abord supposer que tous les points maximaux de $X$ sont des
générisations de $z$, l'ensemble de ces points étant fini; donc l'ensemble des points
maximaux de $U$ est l'ensemble des images $y_i$ par $g$ des points maximaux $x_i$ de
$X$ (aucun point maximal de $X$ ne pouvant par hypothèse être contenu dans $X_y$,
puisque $\dim(\mathcal O_{X,z})\geq2$). Posons $X_i=\overline{\{x_i\}}$. Par
hypothèse, on a $z\in X_i$ et puisque $z$ n'est pas isolé dans $X_y$, il existe dans
$X_y$ un point $x\ne z$. Comme $X_i\ne Z$, il existe dans $X_i$ un point $t_i$ tel
que si l'on pose $T_i=\overline{\{t_i\}}$, on ait $\dim(\mathcal O_{T_i,x})=1$
(10.5.9); cela entraîne par hypothèse $t_i\notin Z$, donc $t_i$ n'est pas une
générisation de $z$. Remplaçant $X$ par $X'=X-\bigcup_iT_i$, on voit que l'image par
$g$ de $X'-X_y$ ne contient pas les images $g(t_i)$, donc ne contient aucune des
composantes irréductibles de $Y'$ dont l'adhérence contient $y$.

Cela étant, comme $X$ et $Y$ sont affines, le morphisme $g:X\to Y$ est affine, donc il
en est de même de la restriction $g':X'\to Y'$; puisque $g'$ est une immersion ouverte,
l'immersion canonique $U\to Y'$ est affine. Posons
$Y_1=\operatorname{Spec}(\mathcal O_{Y,y})$, $Y'_1=Y_1-\{y\}=Y'\cap Y_1$,
$U_1=U\cap Y_1$. Ce qui précède prouve que l'immersion canonique $U_1\to Y'_1$ est
affine, donc $U_1$ est un ouvert affine dans $Y_1$ en vertu de (W), en d'autres termes
l'immersion canonique $U_1\to Y_1$ est affine. Mais puisque cette immersion est de
présentation finie (1.6.2), il résulte de (8.10.5, (viii)) appliqué suivant la méthode de
(8.1.2, a)) qu'en restreignant au besoin $Y$ à un voisinage ouvert affine de $y$, on peut
supposer que l'immersion $U\to Y$ est affine. On en conclurait que l'ouvert $X-X_y$ de
$X$, isomorphe à $U$, serait affine, ce qui contredirait (21.12.7); la proposition est
ainsi démontrée.
\end{proof}

\begin{corollary}[21.12.11]
\label{IV.21.12.11}
Soient $X$, $Y$ deux préschémas localement noethériens, $X$ étant supposé
irréductible; soit $g:X\to Y$ un morphisme localement de type fini, et soit $V$ le plus
grand ouvert de $X$ tel que la restriction $g|V:V\to Y$ soit un isomorphisme local.
Supposons que pour tout point $y\in Y$, l'anneau local $\mathcal O_{Y,y}$ vérifie la
condition (W) (21.12.8). Alors toute composante irréductible de $T=X-V$ est, ou bien
de codimension $\leq1$, ou bien telle que son point générique $z$ soit isolé dans
$g^{-1}(g(z))$.
\end{corollary}

\begin{proof}
Posons $g(z)=y$. La question ne dépendant que de la fibre $g^{-1}(y)$ et de l'anneau
$\mathcal O_{Y,y}$, on peut par changement de base se borner au cas où
$Y=\operatorname{Spec}(\mathcal O_{Y,y})$ et où $X$ est affine
$((\mathbf I, 3.6.5)$ et $(\mathbf I, 4.5.5))$; remplaçant éventuellement $X$ par un
voisinage ouvert de $z$ dans $X$, on peut supposer que $T$ est irréductible; en outre on
peut se borner au cas où $V$ est non vide. Le morphisme $g$ peut donc être supposé
séparé et $y$ fermé dans $Y$; puisque la restriction $g|V:V\to Y$ est un isomorphisme
local et que l'ouvert non vide $V$ de $X$ est irréductible, il résulte de
$(\mathbf I, 8.2.8)$ que $g|V:V\to Y$ est une immersion ouverte. La restriction
$g|(V\cap X_y):V\cap X_y\to\{y\}=\operatorname{Spec}(k(y))$ est donc aussi une
immersion ouverte, ce qui montre que $V\cap X_y$ est, soit vide, soit réduit à un point
$x$ rationnel
\oldpage[IV-4]{310}
sur $k(y)$, donc fermé dans $X_y$ $(\mathbf I, 6.4.2)$. Remplaçant encore
éventuellement $X$ par un voisinage ouvert plus petit de $z$ dans $X$, on peut donc se
borner au cas où $V\cap X_y=\varnothing$, autrement dit $T\supset X_y$; comme d'autre
part $z\in X_y$ est le point générique de $T$, on a
$g(T)\subset\overline{\{y\}}=\{y\}$, autrement dit $T=X_y$. On est alors exactement
dans les conditions d'application de (21.12.10), d'où la conclusion.
\end{proof}

\begin{theorem}[21.12.12]
\label{IV.21.12.12}
\textup{(van der Waerden)}. ---
Soient $X$, $Y$ deux préschémas localement noethériens intègres, $g:X\to Y$ un
morphisme birationnel et localement de type fini. On suppose en outre que $Y$ est
normal et que pour tout $y\in Y$, l'anneau $\mathcal O_{Y,y}$ vérifie la condition (W)
de (21.12.8) (conditions remplies en particulier lorsque $Y$ est localement factoriel
(21.12.9)). Si $V$ est le plus grand ouvert de $X$ tel que la restriction
$g|V:V\to Y$ soit un isomorphisme local, toutes les composantes irréductibles de
$T=X-V$ sont de codimension $1$ dans $X$.
\end{theorem}

\begin{proof}
On notera que puisque $g$ est birationnel, l'ensemble ouvert $V$ n'est pas vide; il
suffit donc de prouver qu'un point maximal $z$ de $T$ ne peut être isolé dans $X_y$,
où $y=g(z)$. Mais, quitte à se restreindre à un voisinage ouvert de $z$, on peut se
borner au cas où $T=X_y$; alors toutes les fibres $g^{-1}(y')$ ($y'\in Y$) seraient
vides ou réduites à un point et il résulterait du « Main theorem » (8.12.10) que $g$
serait un isomorphisme local, contrairement à l'hypothèse $z\in T$.
\end{proof}

\begin{corollary}[21.12.13]
\label{IV.21.12.13}
Supposons vérifiées les hypothèses de (21.12.12) et en outre, supposons que $g$ soit
quasi-fini en chacun des points de $X^{(1)}$ (on rappelle que ce sont les points où
$\dim(\mathcal O_{X,x})=1$); alors $g$ est un isomorphisme local, et si de plus $g$ est
séparé, $g$ est une immersion ouverte.
\end{corollary}

\begin{proof}
Il suffit de prouver la première assertion, la seconde étant une conséquence de la
première et de $(\mathbf I, 8.2.8)$. Tout revient à prouver, avec les notations de
(21.12.12), que $T=\varnothing$; dans le cas contraire, un point générique $z$ d'une
composante irréductible de $T$ appartiendrait à $X^{(1)}$ en vertu de (21.12.12), et par
hypothèse il serait isolé dans $X_y$ avec $y=g(z)$; mais on a vu dans la démonstration
de (21.12.12) que cela n'est pas possible.
\end{proof}

\begin{remarks}[21.12.14]
\label{IV.21.12.14}
(i) La conclusion de (21.12.12) s'applique lorsque $Y$ est régulier et intègre et $X$
intègre, car en vertu du théorème d'Auslander--Buchsbaum (21.11.1), $Y$ est alors
localement factoriel. Par contre, on connaît des exemples où $X$ et $Y$ sont des schémas
algébriques normaux de dimension $3$, sur un corps algébriquement clos de
caractéristique quelconque, et où la conclusion de (21.12.12) n'est plus valable.

(ii) L'ensemble $T$ de (21.12.12) est aussi l'ensemble des points où $g$ est
\emph{ramifié}. En effet, si $g$ est non ramifié en un point $x\in X$, il est aussi non
ramifié dans un voisinage ouvert affine $U$ de $x$ (17.3.7), donc $g|U$ est un morphisme
séparé, quasi-fini et birationnel (17.4.3). Puisque $Y$ est normal, on conclut du « Main
theorem » $(\mathbf{III}, 4.4.9)$ que $g$ est un isomorphisme local au point $x$, donc
$x\in X-T$. Réciproquement, $g$ est évidemment non ramifié en tout point où il est un
isomorphisme local. Ceci justifie le titre donné à cette section.

\oldpage[IV-4]{311}
(iii) Sans supposer que $Y$ soit localement factoriel, mais en supposant par contre que
$X$ est intersection complète relativement à $Y$ (chap.~V), nous verrons au chap.~V
qu'on a un résultat plus précis que (21.12.12), en exprimant $T$ comme cycle
1-codimensionnel d'une section d'un $\mathcal O_X$-Module inversible, canoniquement
associée à $g$. Ceci s'appliquera notamment quand $X$ et $Y$ sont tous deux réguliers,
ou lorsque $g$ est un $S$-morphisme, $X$ et $Y$ étant des préschémas lisses sur $S$.

(iv) On peut se demander si (21.12.10) admet une réciproque; autrement dit, pour un
anneau local noethérien donné $A$, si l'on pose $Y=\operatorname{Spec}(A)$, $y$
désignant le point fermé de $Y$, et si, pour tout préschéma localement noethérien $X$ et
tout morphisme $g:X\to Y$ tel que $g':X'\to Y'$ soit une immersion ouverte, toute
composante irréductible de $X_y$ est de codimension $\leq1$ dans $X$ ou est réduite à
un point, alors est-il vrai que $A$ vérifie la condition (W) de (21.12.8)?

(v) Soient $Y$ un préschéma intègre \emph{régulier}, $X$ un préschéma localement
noethérien \emph{normal}, $g:X\to Y$ un morphisme localement de type fini; supposons
en outre que, si $\eta$ est le point générique de $Y$, la fibre $X_\eta$ soit \emph{étale}
sur $k(\eta)$, et soit $T'$ le complémentaire du plus grand ouvert $V'$ de $X$ tel que
$g|V':V'\to Y$ soit étale; est-il vrai alors que toutes les composantes irréductibles de
$T'$ soient de codimension $1$? Il en est bien ainsi lorsqu'on suppose en outre que $g$
est \emph{localement quasi-fini} (« théorème de pureté » de Zariski--Nagata). On peut
montrer que la conjecture précédente serait conséquence de la suivante (et lui serait
même équivalente si la conjecture de (iv) était vérifiée): pour un anneau local régulier
$A$ contenu dans un anneau local intègre et intégralement clos $B$, tel que $B$ soit un
$A$-module fini, l'ouvert $U$ des points de $\operatorname{Spec}(B)$ en lesquels
$\operatorname{Spec}(B)$ est non ramifié sur $\operatorname{Spec}(A)$ (ou, ce qui
revient au même par (18.10.1), étale sur $\operatorname{Spec}(A)$) est affine.
\end{remarks}

\begin{proposition}[21.12.15]
\label{IV.21.12.15}
Soient $S$ un préschéma, $g:X\to S$ un morphisme plat et localement de présentation
finie, dont les fibres $X_s=g^{-1}(s)$ sont géométriquement irréductibles (4.5.2),
$h:Y\to S$ un morphisme lisse; on note $Y_s=h^{-1}(s)$ les fibres de ce morphisme.
Soit $f:X\to Y$ un $S$-morphisme propre tel que, pour tout point maximal $\eta$ de
$S$, le morphisme $f_\eta:X_\eta\to Y_\eta$ soit un isomorphisme. Alors $f$ est un
isomorphisme.
\end{proposition}

\begin{proof}
Puisque $h$ est plat, tout point maximal de $Y$ est au-dessus d'un point maximal de $S$
(2.3.4), donc la réunion des $Y_\eta$, lorsque $\eta$ parcourt l'ensemble des points
maximaux de $S$, est dense dans $Y$; comme les $f_\eta$ sont des isomorphismes, $f$
est dominant et par suite \emph{surjectif} puisqu'il est propre. Il suffit donc de montrer
que $f$ est une \emph{immersion ouverte}. Compte tenu de (17.9.5), il suffit de prouver
que pour tout $s\in S$, $f_s:X_s\to Y_s$ est une immersion ouverte.

Comme tout $s\in S$ est générisation d'un point maximal $\eta$, on peut déjà, en
faisant le changement de base $\operatorname{Spec}(\mathcal O_{S',s})\to S$, où $S'$
est le sous-préschéma réduit de $S$ ayant $\overline{\{\eta\}}$ comme espace
sous-jacent, se ramener au cas où $S$ est un schéma local et intègre de point fermé $s$:
les fibres aux points $\eta$ et $s$ ne sont en effet pas changées et les propriétés des
morphismes $g$, $h$ et $f$ se conservent par changement de base. En outre, la question
est locale sur $Y$; remplaçant $Y$ par un voisinage ouvert affine $U$ dans $Y$ d'un point
de $Y_s$, et $X$ par $f^{-1}(U)$, qui est quasi-compact puisque $f$ est propre, on voit
qu'on peut supposer en outre que $X$ et $Y$ sont de présentation finie sur $S$.

Il existe alors un sous-anneau local noethérien $A_0$ de
$A=\mathcal O_{S,s}$ tel que $A_0\to A$ soit un homomorphisme local, deux morphismes
de type fini $g_0:X_0\to S_0=\operatorname{Spec}(A_0)$,
$h_0:Y_0\to S_0$ et un $S_0$-morphisme
\oldpage[IV-4]{312}
$f_0:X_0\to Y_0$ tels que $g$, $h$ et $f$ se déduisent de $g_0$, $h_0$, $f_0$ par le
changement de base $S\to S_0$ ((8.9.1) et (5.13.3)); on peut en outre supposer $g_0$
plat (11.2.7), $h_0$ lisse (17.7.9) et $f_0$ propre (8.10.5, (xii)). D'autre part, la
projection sur $S_0$ du point générique $\eta$ de $S$ est le point générique $\eta_0$ de
$S_0$ et il résulte de (2.7.1, (viii)) que le morphisme
$(f_0)_{\eta_0}:(X_0)_{\eta_0}\to(Y_0)_{\eta_0}$ est un isomorphisme; enfin le point
fermé $s$ de $S$ est au-dessus du point fermé $s_0$ de $S_0$, donc la fibre
$(X_0)_{s_0}$ de $g_0$ est géométriquement irréductible (4.5.6).

On voit donc qu'on peut se borner au cas où $S$ est le spectre d'un anneau local intègre
noethérien, de point fermé $s$, affaiblir l'hypothèse sur les fibres en supposant seulement
que $X_s$ et $X_\eta$ sont géométriquement irréductibles, et prouver que sous ces
hypothèses $f_s$ est une immersion ouverte. Il y a alors un anneau de valuation discrète
$A'$ et un morphisme $S'=\operatorname{Spec}(A')\to S$ qui transforme le point fermé
$a$ de $S'$ en $s$ et le point générique $b$ de $S'$ en $\eta$ $(\mathbf{II}, 7.1.9)$;
soient $g':X'\to S'$, $h':Y'\to S'$, $f':X'\to Y'$ les morphismes déduits de $g$,
$h$, $f$ par le changement de base $S'\to S$, qui vérifient de nouveau les hypothèses
vérifiées par $g$, $h$, et $f$ dans l'énoncé (21.12.15); si l'on prouve que
$f'_a:X'_a\to Y'_a$ est une immersion ouverte, il résultera de (2.7.1, (x)) qu'il en
sera de même de $f_s:X_s\to Y_s$. On peut donc finalement se borner au cas où $S$ est
le \emph{spectre d'un anneau de valuation discrète}.

Alors, puisque $h:Y\to S$ est lisse, $Y$ est régulier (17.5.8), et puisque la question
est locale sur $Y$, on peut se borner au cas où $Y$ est intègre. Comme $g:X\to S$ est
plat, les points maximaux de $X$ sont contenus dans $X_\eta$ (2.3.4) et puisque
$f_\eta$ est un isomorphisme, $X$ est irréductible et l'anneau local en son point
générique est un corps; de plus, par platitude (3.3.2), on voit que $X$ n'a pas de cycles
premiers immergés, donc $X$ est réduit (3.2.1) et par suite intègre et $f$ est un
morphisme birationnel et séparé. Pour prouver que $f$ est une immersion ouverte, on peut
donc appliquer le critère (21.12.13); pour voir que $f$ est quasi-fini aux points de
$X^{(1)}$, on remarque que l'assertion est évidente en ceux de ces points qui
appartiennent à $X_\eta$; le seul point de $X^{(1)}$ appartenant à $X_s$ est le point
générique $x$ de $X_s$, en vertu de (6.1.1). Or, il résulte de (2.4.6) et de (14.2.2)
que les morphismes $g$ et $h$ sont équidimensionnels; comme $X_\eta$ et $Y_\eta$ sont
isomorphes, les composantes irréductibles de $X_s$ et de $Y_s$ ont toutes même
dimension. Mais par hypothèse $X_s$ est irréductible et on a vu que $f$ est surjectif,
donc il en est de même de $f_s:X_s\to Y_s$, ce qui entraîne que $Y_s$ est aussi
irréductible; si $y$ est le point générique de $Y_s$, on a donc $f(x)=y$, et la relation
$\dim(X_s)=\dim(Y_s)$ entraîne alors, en vertu de (5.6.6), que
$\dim(f_s^{-1}(y))=0$, autrement dit $f_s$ (et par suite aussi $f$) est bien quasi-fini
au point $x$, ce qui achève la démonstration.
\end{proof}

\begin{corollary}[21.12.16]
\label{IV.21.12.16}
Soient $g:X\to S$ un morphisme propre, plat et de présentation finie, $h:Y\to S$ un
morphisme lisse, propre et de présentation finie, $f:X\to Y$ un $S$-morphisme. On
suppose que les fibres $X_s=g^{-1}(s)$ de $g$ sont géométriquement irréductibles. Soit
$U$ l'ensemble des $s\in S$ tels que $f_s:X_s\to Y_s$ soit un isomorphisme. Alors
$U$ est ouvert et fermé dans $S$, et le morphisme $g^{-1}(U)\to h^{-1}(U)$,
restriction de $f$, est un isomorphisme.
\end{corollary}

\begin{proof}
La dernière assertion résultera de la première et de (17.9.5). On sait déjà (9.6.1,
(xi)) que $U$ est localement constructible dans $S$. En vertu de (1.10.1), il suffit de
prouver que $U$ est stable par \emph{spécialisation} et par \emph{générisation}. Pour
démontrer ces assertions on
\oldpage[IV-4]{313}
peut, par un changement de base $\operatorname{Spec}(\mathcal O_{S,s})\to S$, se
ramener au cas où $S$ est un schéma local de point fermé $s$ et de point générique
$\eta$, et il faut montrer que, pour que $f_s$ soit un isomorphisme, il faut et il suffit
que $f_\eta$ le soit. Or, la suffisance de cette condition résulte de (21.12.15). Pour
montrer qu'elle est nécessaire, on raisonne comme dans la démonstration de (21.12.15)
(en remarquant, avec les mêmes notations, que la projection du point fermé de $S$ est le
point fermé de $S_0$) pour se ramener au cas où en outre $S$ est \emph{noethérien};
mais alors la conclusion résulte de $(\mathbf{III}, 4.6.7$, (ii)).
\end{proof}

\begin{remarks}[21.12.17]
\label{IV.21.12.17}
(i) La conclusion de la proposition (21.12.15) n'est plus valable si on élimine
l'hypothèse que les fibres $X_s$ sont irréductibles. Prenons en effet
$S=\operatorname{Spec}(A)$, où $A$ est un anneau de valuation discrète,
$Y=\mathbb P^1_S$, qui est propre et lisse sur $S$ (17.3.9). Notons encore $s$ et
$\eta$ le point fermé et le point générique de $S$; soit $z$ un point fermé de $Y_s$,
par exemple un point rationnel sur $k=k(s)$, et soit $X$ le $Y$-préschéma obtenu en
faisant éclater le point $z$. Comme l'anneau de polynômes $A[T]$ est régulier
$(\mathbf 0, 17.3.7)$ et de dimension $2$, on voit comme dans la démonstration de
(15.1.1.6) que, si $f:X\to Y$ est le morphisme structural, $f^{-1}(z)$ est isomorphe
à $\mathbb P^1_k$, et d'autre part, le complémentaire de $f^{-1}(z)$ dans $X_s$ est
isomorphe à $Y_s-\{z\}$, donc $Z_1=f^{-1}(z)$ et $Z_2$, adhérence dans $X_s$ du
complémentaire de $Z_1$, sont les deux composantes irréductibles de $X_s$. Par
ailleurs, $f$ est évidemment propre et $g=h\circ f$ est plat, car $X$ est intègre
$(\mathbf{II}, 8.1.4)$ et pour tout ouvert affine $U$ de $X$, l'homomorphisme
$A\to\Gamma(U,\mathcal O_X)$ est injectif $(\mathbf I, 1.2.7)$, donc
$\Gamma(U,\mathcal O_X)$ est un $A$-module sans torsion, et par suite plat puisque $A$
est un anneau de valuation discrète $(\mathbf 0_{\mathrm I}, 6.3.4)$. Il est clair que
$f_\eta:X_\eta\to Y_\eta$ est un isomorphisme, bien que
$f_s:X_s\to Y_s$ n'en soit pas un.

(ii) La conclusion de (21.12.15) tombe aussi en défaut si, dans les hypothèses, on
supprime l'hypothèse que $f$ est \emph{propre}. Il suffit pour le voir de définir $S$ et
$Y$ comme dans (i), mais de remplacer $X$ par le complémentaire $X'$ de $Z_2$ dans le
préschéma $X$ défini dans (i), $f$ par la restriction $f':X'\to Y$ de $f$; le morphisme
$g':X'\to S$, restriction de $g$, est encore plat, et cette fois $X'_s$ est
géométriquement irréductible; $X'_s$ est d'ailleurs isomorphe au complémentaire d'un
point fermé dans $\mathbb P^1_k$, donc est un schéma affine isomorphe à $\mathbb A^1_k$;
comme son image dans $Y_s$ est réduite au point fermé $z$, $f'$ n'est pas propre
$(\mathbf{III}, 4.4.2)$ et $f'_s$ n'est pas un isomorphisme, bien que $f'_\eta$ en soit
un.

(iii) Il se peut que l'énoncé de la proposition (21.12.15) reste valable quand on y
remplace le mot « isomorphisme » par « morphisme étale » (cf. (21.12.14, (v))). Il en
sera alors encore de même de (21.12.16).
\end{remarks}
\subsection{Couples parafactoriels. Anneaux locaux parafactoriels}
\label{subsection:IV.21.13}
\label{IV.21.13}

\begin{definition}[21.13.1]
\label{IV.21.13.1}
Soient $X$ un espace annelé, $Y$ une partie fermée de $X$, $U=X-Y$.
On dit que le couple $(X,Y)$ est \emph{parafactoriel} si, pour tout ouvert $V$ de $X$,
le foncteur restriction
$\mathscr L\mapsto\mathscr L|(U\cap V)$, de la catégorie des
$\mathcal O_V$-Modules inversibles dans celle des $\mathcal O_{U\cap V}$-Modules
inversibles, est une \emph{équivalence de catégories}.
\end{definition}

Dire que le couple $(X,Y)$ est parafactoriel signifie donc que, pour tout ouvert $V$
de $X$, les conditions suivantes sont vérifiées :
\begin{enumerate}[label=\arabic*\textsuperscript{o}]
  \item le foncteur $\mathscr L\mapsto\mathscr L|(U\cap V)$ est \emph{pleinement
  fidèle}, autrement dit, pour deux $\mathcal O_V$-Modules inversibles $\mathscr L$,
  $\mathscr L'$, l'application de restriction
  \[
    \operatorname{Hom}_{\mathcal O_V}(\mathscr L,\mathscr L')
    \longrightarrow
    \operatorname{Hom}_{\mathcal O_{U\cap V}}
    (\mathscr L|(U\cap V),\mathscr L'|(U\cap V))
  \]
  est \emph{bijective};
  \item le foncteur $\mathscr L\mapsto\mathscr L|(U\cap V)$ est
  \emph{essentiellement surjectif}, autrement dit, pour tout
  $\mathcal O_{U\cap V}$-Module inversible $\mathscr L_0$, il existe un
  $\mathcal O_V$-Module inversible $\mathscr L$ tel que $\mathscr L_0$ soit
  \oldpage[IV-4]{314}
  isomorphe à $\mathscr L|(U\cap V)$; cela s'exprime encore en disant que
  l'homomorphisme canonique (21.3.2.4)
  \[
    \operatorname{Pic}(V)\longrightarrow\operatorname{Pic}(U\cap V)
  \]
  est \emph{surjectif}.
\end{enumerate}

\begin{lemma}[21.13.2]
\label{IV.21.13.2}
Soit $f:X'\to X$ un morphisme d'espaces annelés; pour tout ouvert $V$ de $X$, on
note $f_V:f^{-1}(V)\to V$ la restriction de $f$. Les conditions suivantes sont
équivalentes :
\begin{enumerate}[label=\emph{\alph*)}]
  \item Pour tout ouvert $V$ de $X$, le foncteur
  $\mathscr E\mapsto f_V^*(\mathscr E)$ de la catégorie des $\mathcal O_V$-Modules
  localement libres de rang fini dans la catégorie des
  $\mathcal O_{f^{-1}(V)}$-Modules localement libres de rang fini, est fidèle
  (resp.\ pleinement fidèle).
  \item[\emph{a')}] Pour tout ouvert $V$ de $X$, le foncteur
  $\mathscr L\mapsto f_V^*(\mathscr L)$ de la catégorie des $\mathcal O_V$-Modules
  localement libres de rang $1$ dans la catégorie des
  $\mathcal O_{f^{-1}(V)}$-Modules localement libres de rang $1$ est fidèle
  (resp.\ pleinement fidèle).
  \item Pour tout ouvert $V$ de $X$, l'homomorphisme canonique
  $(\mathbf 0_{\mathrm I}, 4.4.3.2)$
  $\mathscr E\to(f_V)_*(f_V^*(\mathscr E))$ est un monomorphisme
  (resp.\ un isomorphisme) pour tout $\mathcal O_V$-Module localement libre
  $\mathscr E$ de rang fini.
  \item[\emph{b')}] L'homomorphisme canonique
  $\mathcal O_X\to f_*(\mathcal O_{X'})$ est un monomorphisme
  (resp.\ un isomorphisme).
\end{enumerate}

Supposons que l'homomorphisme canonique $\mathcal O_X\to f_*(\mathcal O_{X'})$ soit
bijectif; alors, pour qu'un $\mathcal O_{X'}$-Module localement libre de rang fini
$\mathscr E'$ soit isomorphe à un $\mathcal O_{X'}$-Module de la forme
$f^*(\mathscr E)$, où $\mathscr E$ est un $\mathcal O_X$-Module localement libre de
rang fini, il faut et il suffit que les deux conditions suivantes soient remplies :
\begin{enumerate}[label=(\roman*)]
  \item $f_*(\mathscr E')$ est un $\mathcal O_X$-Module localement libre;
  \item l'homomorphisme canonique $(\mathbf 0_{\mathrm I}, 4.4.3.3)$
  $f^*(f_*(\mathscr E'))\to\mathscr E'$ est un isomorphisme.
\end{enumerate}
Lorsque ces deux conditions sont remplies, $\mathscr E$ est isomorphe à
$f_*(\mathscr E')$.
\end{lemma}

\begin{proof}
Dire que le foncteur $\mathscr E\mapsto f_V^*(\mathscr E)$ est fidèle
(resp.\ pleinement fidèle) signifie que pour deux $\mathcal O_V$-Modules localement
libres de rang fini $\mathscr E_1$, $\mathscr E_2$, l'application
\[
  \operatorname{Hom}_{\mathcal O_V}(\mathscr E_1,\mathscr E_2)
  \longrightarrow
  \operatorname{Hom}_{\mathcal O_{f^{-1}(V)}}
  (f_V^*(\mathscr E_1),f_V^*(\mathscr E_2))
\]
est injective (resp.\ bijective); comme ceci doit avoir lieu en remplaçant $V$ par un
ouvert $W\subset V$ et $\mathscr E_1$, $\mathscr E_2$ par $\mathscr E_1|W$,
$\mathscr E_2|W$, et que
\[
  \operatorname{Hom}_{\mathcal O_W}(\mathscr E_1|W,\mathscr E_2|W)
  =\Gamma(W,\mathscr Hom_{\mathcal O_V}(\mathscr E_1,\mathscr E_2)),
\]
la condition signifie encore que l'homomorphisme canonique de faisceaux
\[
\label{IV.21.13.2.1}
  \mathscr Hom_{\mathcal O_V}(\mathscr E_1,\mathscr E_2)
  \longrightarrow
  (f_V)_*(f_V^*(\mathscr Hom_{\mathcal O_V}(\mathscr E_1,\mathscr E_2)))
  \tag{21.13.2.1}
\]
est injectif (resp.\ bijectif). Mais puisque $\mathscr E_1$ et $\mathscr E_2$ sont
localement libres de rang fini,
$\mathscr Hom_{\mathcal O_V}(\mathscr E_1,\mathscr E_2)$, isomorphe à
$\mathscr E_1^\vee\otimes_{\mathcal O_V}\mathscr E_2$
$(\mathbf 0_{\mathrm I}, 5.4.2)$, est aussi localement libre de rang fini; cela montre
déjà que $b)$ entraîne $a)$, et inversement $b)$ est un cas particulier de $a)$,
compte tenu de l'isomorphisme
$\mathscr E\xrightarrow{\sim}\mathscr Hom_{\mathcal O_V}(\mathcal O_V,\mathscr E)$.
Il est clair que $a')$ est un cas particulier de $a)$ et $b')$ un cas particulier de
$a')$. Inversement, comme $b)$ est une propriété locale, on peut, pour la vérifier,
se borner au cas où $\mathscr E=\mathcal O_V^n$, et cela montre que $b')$ entraîne
$b)$.

Si l'homomorphisme canonique $\mathcal O_X\to f_*(\mathcal O_{X'})$ est bijectif,
et si les conditions (i) et (ii) sont remplies, il est clair que
$\mathscr E'=f^*(\mathscr E)$ avec $\mathscr E=f_*(\mathscr E')$, à un isomorphisme
près.
\oldpage[IV-4]{315}
Inversement, supposons que $\mathscr E'=f^*(\mathscr E)$, avec $\mathscr E$ localement
libre; la question étant locale sur $X$, on peut supposer que
$\mathscr E=\mathcal O_X^n$, d'où $\mathscr E'=\mathcal O_{X'}^n$, et les conditions
(i) et (ii) résultent de l'hypothèse que
$\mathcal O_X\to f_*(\mathcal O_{X'})$ est un isomorphisme.
\end{proof}

Dans ce numéro, nous appliquerons le lemme précédent à une injection canonique
$j:U\to X$, où $U$ est l'espace annelé induit par $X$ sur un \emph{ouvert} de $X$.
Avec ces notations, (21.13.2) se traduit en :

\begin{corollary}[21.13.3]
\label{IV.21.13.3}
Soient $X$ un espace annelé, $Y$ une partie fermée de $X$, $U=X-Y$,
$j:U\to X$ l'injection canonique. Les conditions suivantes sont équivalentes :
\begin{enumerate}[label=\emph{\alph*)}]
  \item Pour tout ouvert $V$ de $X$, le foncteur restriction
  $\mathscr E\mapsto\mathscr E|(U\cap V)$ de la catégorie des
  $\mathcal O_V$-Modules localement libres de rang fini dans la catégorie des
  $\mathcal O_{U\cap V}$-Modules localement libres de rang fini est fidèle
  (resp.\ pleinement fidèle).
  \item[\emph{a')}] Pour tout ouvert $V$ de $X$, le foncteur restriction
  $\mathscr L\mapsto\mathscr L|(U\cap V)$ de la catégorie des
  $\mathcal O_V$-Modules localement libres de rang $1$ dans la catégorie des
  $\mathcal O_{U\cap V}$-Modules localement libres de rang $1$ est fidèle
  (resp.\ pleinement fidèle).
  \item Pour tout ouvert $V$ de $X$, l'homomorphisme canonique
  $\mathscr E\to(j_V)_*(\mathscr E|(U\cap V))$ est injectif (resp.\ bijectif) pour
  tout $\mathcal O_V$-Module localement libre de rang fini $\mathscr E$.
  \item[\emph{b')}] L'homomorphisme canonique
  $\mathcal O_X\to j_*(\mathcal O_U)$ est injectif (resp.\ bijectif).
\end{enumerate}

Supposons que l'homomorphisme canonique $\mathcal O_X\to j_*(\mathcal O_U)$ soit
bijectif. Alors, pour qu'un $\mathcal O_U$-Module localement libre de rang fini
$\mathscr F$ soit de la forme $\mathscr E|U$, où $\mathscr E$ est un
$\mathcal O_X$-Module localement libre de rang fini, il faut et il suffit que
$j_*(\mathscr F)$ soit un $\mathcal O_X$-Module localement libre, et dans ce cas, on
peut prendre $\mathscr E=j_*(\mathscr F)$.
\end{corollary}

\begin{lemma}[21.13.4]
\label{IV.21.13.4}
Soient $X$ un préschéma localement noethérien, $Y$ une partie fermée de $X$,
$U=X-Y$, $j:U\to X$ l'injection canonique. Pour que l'homomorphisme canonique
$\mathcal O_X\to j_*(\mathcal O_U)$ soit injectif (resp.\ bijectif), il faut et il
suffit que $\operatorname{prof}_Y(\mathcal O_X)\geq1$
(resp.\ $\operatorname{prof}_Y(\mathcal O_X)\geq2$) (autrement dit, que
$\operatorname{prof}(\mathcal O_{X,x})\geq1$
(resp.\ $\operatorname{prof}(\mathcal O_{X,x})\geq2$) pour tout $x\in Y$).
\end{lemma}

\begin{proof}
Ces assertions sont en effet des cas particuliers de (5.10.2) et (5.10.5).
\end{proof}

\begin{proposition}[21.13.5]
\label{IV.21.13.5}
Soient $X$ un espace annelé, $Y$ une partie fermée de $X$, $U=X-Y$,
$j:U\to X$ l'injection canonique. Pour que le couple $(X,Y)$ soit parafactoriel, il
faut et il suffit que l'homomorphisme canonique $\mathcal O_X\to j_*(\mathcal O_U)$
soit bijectif, et que pour tout ouvert $V$ de $X$ et tout
$\mathcal O_{U\cap V}$-Module inversible $\mathscr L$, $(j_V)_*(\mathscr L)$ soit un
$\mathcal O_V$-Module inversible (notation de (21.13.3)).
\end{proposition}

\begin{proof}
C'est une conséquence immédiate de la définition (21.13.1) et de (21.13.3).
\end{proof}

\begin{corollary}[21.13.6]
\label{IV.21.13.6}
Soient $X$ un espace annelé, $Y$ une partie fermée de $X$.
\begin{enumerate}[label=(\roman*)]
  \item Si le couple $(X,Y)$ est parafactoriel, il en est de même de
  $(W,Y\cap W)$ pour tout ouvert $W$ de $X$. Inversement, si $(W_\alpha)$ est un
  recouvrement ouvert de $X$ tel que chacun des couples
  $(W_\alpha,Y\cap W_\alpha)$ soit parafactoriel, alors le couple $(X,Y)$ est
  parafactoriel.
  \item Si le couple $(X,Y)$ est parafactoriel, il en est de même de $(X,Y')$ pour
  toute partie fermée $Y'$ de $Y$.
  \item Supposons que $X$ soit un préschéma, et soient $X'$ un préschéma,
  $f:X'\to X$ un morphisme fidèlement plat et quasi-compact; posons
  $Y'=f^{-1}(Y)$. Supposons que le couple $(X',Y')$ soit parafactoriel et que
  l'ouvert $U=X-Y$ soit rétrocompact dans $X$ $(\mathbf 0_{\mathrm{III}}, 9.1.1)$.
  Alors le couple $(X,Y)$ est parafactoriel.
\end{enumerate}
\end{corollary}

\begin{proof}
\oldpage[IV-4]{316}
(i) Comme le fait que l'homomorphisme canonique
$\mathcal O_X\to j_*(\mathcal O_U)$ soit bijectif est une propriété locale sur $X$,
tout revient à voir (en vertu de (21.13.5)), en désignant par
$j_\alpha:U\cap W_\alpha\to W_\alpha$ l'injection canonique, que l'on a la propriété
suivante : si, pour tout $\alpha$ et pour tout $\mathcal O_U$-Module inversible
$\mathscr L$, $(j_\alpha)_*(\mathscr L|(U\cap W_\alpha))$ est un
$\mathcal O_{W_\alpha}$-Module inversible, alors $j_*(\mathscr L)$ est un
$\mathcal O_X$-Module inversible. Mais la propriété d'être un $\mathcal O_X$-Module
inversible est locale sur $X$, et $(W_\alpha)$ est un recouvrement ouvert de $X$;
comme on a
\[
  j_*(\mathscr L)|W_\alpha
  =(j_\alpha)_*(\mathscr L|(U\cap W_\alpha)),
\]
cela prouve notre assertion.

(ii) Posons $U'=X-Y'$, et soit $j':U'\to X$ l'injection canonique. Dire que
l'homomorphisme $\mathcal O_X\to j'_*(\mathcal O_{U'})$ est bijectif revient à dire
que pour tout ouvert $V$ de $X$, l'homomorphisme
$\Gamma(V,\mathcal O_X)\to\Gamma(V\cap U',\mathcal O_X)$ est bijectif; mais
l'homomorphisme composé
\[
  \Gamma(V,\mathcal O_X)\longrightarrow\Gamma(V\cap U',\mathcal O_X)
  \longrightarrow\Gamma(V\cap U,\mathcal O_X)
\]
est bijectif par hypothèse (21.13.5), et en remplaçant $V$ par $V\cap U'$, on voit que
$\Gamma(V\cap U',\mathcal O_X)\to\Gamma(V\cap U,\mathcal O_X)$ est bijectif, donc
$\Gamma(V,\mathcal O_X)\to\Gamma(V\cap U',\mathcal O_X)$ est bijectif. Notons ensuite
que si $j'':U\to U'$ est l'injection canonique et si $\mathscr L'$ est un
$\mathcal O_{U'}$-Module inversible, $\mathscr L'|U$ est un $\mathcal O_U$-Module
inversible, donc
$j_*(\mathscr L'|U)=j'_*(j''_*(\mathscr L'|U))$ est par hypothèse un
$\mathcal O_X$-Module inversible. Comme le couple $(U',U'\cap Y)$ est parafactoriel
en vertu de (i), on a $j''_*(\mathscr L'|U)\cong\mathscr L'$, donc
$j'_*(\mathscr L')$ est un $\mathcal O_X$-Module inversible. Il suffit alors pour
conclure de remplacer dans le raisonnement précédent $X$, $U$ et $U'$ par $V$,
$V\cap U$ et $V\cap U'$, où $V$ est un ouvert de $X$.

(iii) Posons $U'=X'-Y'=f^{-1}(U)$ et notons que puisqu'il s'agit de préschémas, on
peut écrire $U'=U\times_X X'$; soit $f_U:f^{-1}(U)\to U$ la restriction de $f$ et
soit $j':U'\to X'$ l'injection canonique, qui s'écrit aussi $j_{(X')}$. Montrons
d'abord que l'homomorphisme canonique
$\rho:\mathcal O_X\to j_*(\mathcal O_U)$ est bijectif; pour cela, notons que par
hypothèse l'homomorphisme canonique
$\mathcal O_{X'}\to j'_*(\mathcal O_{U'})$ est bijectif; mais puisque le morphisme
$j$ est quasi-compact et séparé par hypothèse, et le morphisme $f$ plat,
$j'_*(\mathcal O_{U'})=j'_*(f_U^*(\mathcal O_U))$ s'identifie canoniquement à
$f^*(j_*(\mathcal O_U))$ en vertu de (2.3.1); comme
$\mathcal O_{X'}=f^*(\mathcal O_X)$, on voit que l'homomorphisme
$f^*(\rho):f^*(\mathcal O_X)\to f^*(j_*(\mathcal O_U))$ est bijectif; on en conclut
que $\rho$ est bijectif puisque $f$ est fidèlement plat (2.2.7). Considérons ensuite
un $\mathcal O_U$-Module inversible $\mathscr L$, et soit
$\mathscr L'=f_U^*(\mathscr L)$, qui est un $\mathcal O_{U'}$-Module inversible. Par
hypothèse $j'_*(f_U^*(\mathscr L))$ est un $\mathcal O_{X'}$-Module inversible, qui
est isomorphe à $f^*(j_*(\mathscr L))$ par (2.3.1) comme ci-dessus. Mais puisque $f$
est fidèlement plat et quasi-compact, cela entraîne que $j_*(\mathscr L)$ est un
$\mathcal O_X$-Module localement libre (2.5.2). Pour achever la démonstration, il
suffit de remplacer $X$ par un ouvert $V$ et $X'$ par $f^{-1}(V)$ dans ce qui
précède.
\end{proof}

\begin{definition}[21.13.7]
\label{IV.21.13.7}
On dit qu'un anneau local $A$ est \emph{parafactoriel} si, posant
$X=\operatorname{Spec}(A)$ et désignant par $a$ le point fermé de $X$, le couple
$(X,\{a\})$ est parafactoriel.
\end{definition}

\begin{proposition}[21.13.8]
\label{IV.21.13.8}
Les notations étant celles de (21.13.7), on pose $U=X-\{a\}$. Pour que $A$ soit
parafactoriel, il faut et il suffit qu'il vérifie les conditions suivantes :
\begin{enumerate}[label=(\roman*)]
  \item L'homomorphisme canonique
  $A=\Gamma(X,\mathcal O_X)\to\Gamma(U,\mathcal O_X)$ est bijectif.
  \item $\operatorname{Pic}(U)=0$.
\end{enumerate}
\oldpage[IV-4]{317}
Si en outre $A$ est noethérien, la condition (i) équivaut à
\begin{enumerate}
  \item[\textup{(i bis)}] $\operatorname{prof}(A)\geq2$.
\end{enumerate}
\end{proposition}

\begin{proof}
En effet, le seul ouvert de $X$ contenant $a$ est $X$ lui-même, et par suite tout
$\mathcal O_X$-Module inversible est isomorphe à $\mathcal O_X$, autrement dit
$\operatorname{Pic}(X)=0$; la première assertion résulte donc de la définition
(21.13.1) et de (21.13.3). La seconde assertion n'est autre qu'un cas particulier de
(21.13.4).
\end{proof}

\begin{env}[21.13.9]
\label{IV.21.13.9}
\emph{Exemples}. ---

(i) Un anneau local noethérien parafactoriel est nécessairement de dimension
$\geq2$ en vertu de (21.13.8); en d'autres termes un anneau local noethérien de
dimension $\leq1$ n'est pas parafactoriel.

(ii) Un anneau local noethérien \emph{factoriel} $A$ de dimension $\geq2$ est
parafactoriel, comme il résulte de (21.13.8) et de (21.6.14).

(iii) Si $A$ est un anneau local noethérien de dimension $\geq3$ et parafactoriel, il
n'est pas nécessairement normal ni même réduit. Considérons un anneau local régulier
$B$ de dimension $\geq3$, et soit $A=D_B(B)$ $(\mathbf 0, 18.2.3)$, isomorphe à
$B[T]/(T^2)$; on voit aussitôt que l'on a
$\operatorname{prof}(A)=\operatorname{prof}(B)=\dim(B)\geq3$. Pour montrer que $A$
est parafactoriel, il suffit donc, avec les notations de (21.13.8), de prouver que
$\operatorname{Pic}(U)=0$. Soit $\mathfrak J$ le noyau de l'augmentation $A\to B$,
qui est tel que $\mathfrak J^2=0$ et qui, en tant que $B$-module, est isomorphe à $B$.
Comme $B=A/\mathfrak J$, $X=\operatorname{Spec}(A)$ et
$X_0=\operatorname{Spec}(B)$ ont même espace sous-jacent; si
$\mathscr J=\widetilde{\mathfrak J}$, $X_0$ est le sous-schéma défini par
$\mathscr J$; nous noterons $U_0$ le sous-préschéma induit par $X_0$ sur l'ouvert
$U$. Pour tout $z\in\mathfrak J$, posons $\varphi(z)=1+z$; puisque
$(1+z)(1-z)=1$, $1+z$ est inversible dans $A$ et $\varphi(\mathfrak J)$ est le noyau
de l'homomorphisme surjectif canonique de groupes multiplicatifs $A^*\to B^*$; en
d'autres termes, on a une suite exacte de groupes commutatifs
\[
  0\longrightarrow\mathfrak J\xrightarrow{\varphi}A^*\longrightarrow B^*
  \longrightarrow1
\]
(les trois derniers groupes étant multiplicatifs, les deux premiers additifs). Pour la
même raison, pour tout $t\in A$, on a la suite exacte
\[
  0\longrightarrow\mathfrak J_t\xrightarrow{\varphi_t}(A_t)^*
  \longrightarrow(B_{t_0})^*\longrightarrow1
\]
en désignant par $t_0$ l'image de $t$ dans $B$, puisque
$B_{t_0}=A_t/\mathfrak J_t$; autrement dit, on a une suite exacte de faisceaux de
groupes commutatifs sur l'espace topologique $X$
\[
  0\longrightarrow\mathscr J\xrightarrow{u}\mathcal O_X^*
  \longrightarrow\mathcal O_{X_0}^*\longrightarrow1
\]
d'où par restriction à l'ouvert $U$, une suite exacte
\[
  0\longrightarrow\mathscr J|U\longrightarrow\mathcal O_U^*
  \longrightarrow\mathcal O_{U_0}^*\longrightarrow1.
\]
Par la suite exacte de cohomologie, on en déduit une suite exacte
\[
\label{IV.21.13.9.1}
  \mathrm H^1(U,\mathscr J)\longrightarrow\mathrm H^1(U,\mathcal O_U^*)
  \xrightarrow{u_1}\mathrm H^1(U,\mathcal O_{U_0}^*)
  \tag{21.13.9.1}
\]
Mais puisque $\mathfrak J$ est un $B$-module isomorphe à $B$, on a
$\mathrm H^1(U,\mathscr J)=\mathrm H^1(U_0,\mathcal O_{U_0})$, et il résulte du
chap.~III, 3\textsuperscript{e} partie (voir aussi [41, III, Exemple III-1]) que l'on a
\oldpage[IV-4]{318}
$\mathrm H^1(U_0,\mathcal O_{U_0})=0$ en raison de la relation
$\operatorname{prof}(B)\geq3$. Par ailleurs, comme $B$ est factoriel, on a
$\mathrm H^1(U,\mathcal O_{U_0}^*)=\operatorname{Pic}(U_0)=0$; on tire donc bien de
la suite exacte précédente que
$\operatorname{Pic}(U)=\mathrm H^1(U,\mathcal O_U^*)=0$.

(iv) Il existe aussi des anneaux locaux noethériens de dimension $3$ qui sont
intègres, intégralement clos et parafactoriels, mais non factoriels. Soit en effet $B$
un anneau local noethérien, complet, intègre et intégralement clos, de dimension
$\geq2$ et non factoriel (par exemple le complété de l'anneau localisé de l'algèbre
intègre $k[U,V,W]/(W^2-UV)$ en l'idéal image de $(U)+(V)+(W)$). Alors il résulte de
ce que l'on verra plus bas (21.14.2) que l'anneau local $A=B[[T]]$ des séries formelles
en $T$ sur $B$ est parafactoriel, mais il n'est pas factoriel, sans quoi $B$ le serait
en vertu de (21.13.12) ci-dessous.

(v) On peut montrer qu'un anneau local d'intersection complète absolue (19.3.1) de
dimension $\geq4$, est parafactoriel (cf. [41, XI, 3.13 (i)]).

(vi) On a vu (Remarque (ii)) que tout anneau local noethérien factoriel $A$ de
dimension $2$ est parafactoriel. Mais il y a des anneaux locaux noethériens de
dimension $2$ qui sont parafactoriels et non factoriels. On peut montrer que, pour
qu'un anneau local noethérien $A$ de dimension $2$ soit parafactoriel et non
factoriel, il faut et il suffit qu'il satisfasse aux trois conditions suivantes :
\begin{enumerate}[label=\arabic*\textsuperscript{o}]
  \item $A$ est un anneau de Cohen--Macaulay (autrement dit
  $\operatorname{prof}(A)=2$).
  \item $A$ est intègre et si $A'$ est sa clôture intégrale, $A'$ est factoriel et est
  une $A$-algèbre finie.
  \item Soit $\mathfrak J$ le conducteur de $A$ dans $A'$ (annulateur du $A$-module
  $A'/A$, ou encore le plus grand idéal de $A'$ contenu dans $A$); posons
  $B=A/\mathfrak J$, $B'=A'/\mathfrak J$; alors $\dim(B)=1$ (ce qui implique
  $A'\neq A$, autrement dit $A$ n'est pas intégralement clos), et l'application
  canonique $\operatorname{Div}(B)\to\operatorname{Div}(B')$ (21.4.5) est
  \emph{surjective}.
\end{enumerate}
On peut montrer en outre que ces conditions entraînent la propriété suivante :

$4^{\mathrm o}$ L'anneau $B'$ (et \emph{a fortiori} $B\subset B'$) est réduit, et le
morphisme $\operatorname{Spec}(B')\to\operatorname{Spec}(B)$ est bijectif.

Si l'on pose $X=\operatorname{Spec}(A)$, $X'=\operatorname{Spec}(A')$,
$Y=\operatorname{Spec}(B)$, $Y'=\operatorname{Spec}(B')$, de sorte que $Y$
(resp.\ $Y'$) est défini par l'idéal $\mathfrak J$ de $A$ (resp.\ $A'$), le morphisme
structural $f:X'\to X$ est un isomorphisme de $X'-Y'$ sur $X-Y$ (Bourbaki,
\emph{Alg. comm.}, chap.~V, \S1, n\textsuperscript{o}~5, cor.~5 de la prop.~16). On
voit donc en vertu de $4^{\mathrm o}$ que $f$ est un morphisme \emph{bijectif}; en
d'autres termes, $X$ est un préschéma \emph{unibranche} (6.15.1) (et en particulier
$A'$ est un anneau local noethérien); en général $X$ n'est pas géométriquement
unibranche. L'espace $Y$, de dimension $1$, est constitué par le point fermé $x$ de
$X$ et les points maximaux $y_i$ ($1\leq i\leq r$) de $Y$, et l'unique point $y_i'$
de $Y'$ au-dessus de $y_i$ est aussi un point maximal de $Y'$; comme $B$ et $B'$ sont
réduits, on en déduit que l'on a
\[
  \mathfrak m_{y_i}\mathcal O_{X',y_i'}=\mathfrak m_{y_i'}.
\]
Pour $z=y_i$ et $z'=y_i'$, cette relation est aussi vérifiée de façon évidente pour
$z\notin Y$ et $z'$ l'unique point de $f^{-1}(z)$, donc pour tout
$z\in U=X-\{x\}$ et $z'$ l'unique point de $f^{-1}(z)$. En particulier, si l'anneau
$A$ est de caractéristique $0$ (0, 21.1.1) on voit que $U'$ est \emph{non ramifié} sur
$U$ (mais non étale en général).

Nous nous bornerons ici à démontrer que les conditions $1^{\mathrm o}$,
$2^{\mathrm o}$, $3^{\mathrm o}$ ci-dessus sont \emph{suffisantes} pour que $A$ soit
parafactoriel. Or, en vertu de la condition $1^{\mathrm o}$ et de (21.13.8), il suffit
de démontrer que $\operatorname{Pic}(U)=0$. Puisque $A'$ est factoriel en vertu de
$2^{\mathrm o}$, on a $\operatorname{Pic}(U')=0$, et on déduit de (21.8.5, (ii)) une
suite exacte
\[
  1\longrightarrow
  \left(\prod_{i=1}^r
  \mathcal O_{U',y_i'}^*/\mathcal O_{U,y_i}^*\right)
  \big/\operatorname{Im}\Gamma(U',\mathcal O_{U'}^*)
  \longrightarrow\operatorname{Pic}(U)\longrightarrow\operatorname{Pic}(U')
  \longrightarrow0
\]
et il s'agit donc de montrer que le second terme de cette suite est réduit à l'élément
neutre, en prenant pour $S$ l'ensemble des $y_i$ ($1\leq i\leq r$) et en notant
qu'ici $f_*(\mathcal O_{X'})$ est la $\mathcal O_X$-Algèbre $\widetilde{A'}$. Soient
$\mathfrak p_i$, $\mathfrak p_i'$ les idéaux premiers de $A$ et $A'$ correspondant
aux points $y_i$, $y_i'$, et remarquons que $\mathfrak p_i'\cap A=\mathfrak p_i$;
donnons-nous pour chaque $i$ un élément inversible $a_i'/s_i$ de
$A'_{\mathfrak p_i'}$, avec $a_i'\in A'$, $s_i\notin\mathfrak p_i'$; comme nous
n'avons à examiner que les groupes quotients
$(A'_{\mathfrak p_i'})^*/(A_{\mathfrak p_i})^*$, on peut supposer $s_i=1$ pour tout
$i$; soit $b_i'\in B'=A'/\mathfrak J$ l'image canonique de $a_i'$, de sorte que
$b_i'/1$ est un élément non nul de $k(y_i')=\mathcal O_{Y',y_i'}$. L'ensemble
$U\cap Y$ formé par les $y_i$ est un ouvert affine de la forme $D(t)$ dans $Y$, avec
$t\in\mathfrak m$, idéal maximal de $A$; il existe donc un élément inversible
$b'\in B'_t$ dont les $b_i'/1$ sont les images canoniques, et comme $t$ est régulier
dans l'anneau intègre $A'$, on peut supposer $b'$ de la forme $b''/t^m$, où
$b''\in B'$ est inversible. Soit $a''$ un élément de $A'$ dans la classe $b''$, qui
est donc nécessairement inversible; $a'=a''/t^m$ est inversible dans $A'_t$ et pour
tout $i$, on a
\[
  u_i=a''-a_i't^m\in\mathfrak J,
  \qquad
  t^m a_i'=a''-u_i=a''(1-a''^{-1}u_i);
\]
mais $a''^{-1}u_i\in\mathfrak J\subset\mathfrak m$, donc $1-a''^{-1}u_i$ est un
élément inversible de $A$, et les classes de $a'$ et $a_i'$ dans
$(A'_{\mathfrak p_i'})^*/(A_{\mathfrak p_i})^*$ sont les mêmes, ce qui achève la
démonstration.

\oldpage[IV-4]{319}
Pour avoir un exemple explicite d'anneau parafactoriel de dimension $2$ obtenu de
cette manière et \emph{non factoriel}, considérons l'anneau
\[
  E=\mathbb R[[U,V]]/(U^2+V^2),
\]
dont la clôture intégrale $E'$ s'identifie à $\mathbb C[[U]]$ (6.15.11). Si l'on pose
$A=E[[T]]$, la clôture intégrale de $A$ est l'anneau $A'=E'[[T]]$ (Bourbaki,
\emph{Alg. comm.}, chap.~V, \S1, n\textsuperscript{o}~4, prop.~14); on vérifie
aussitôt que le conducteur de $E$ dans $E'$ est l'idéal maximal $\mathfrak n$ de
$E$, donc le conducteur $\mathfrak J$ de $A$ dans $A'$ est $\mathfrak n[[T]]$, et on
a $B=A/\mathfrak J\cong\mathbb R[[T]]$, $B'=A'/\mathfrak J\cong\mathbb C[[T]]$.
Il est alors immédiat que les conditions $1^{\mathrm o}$, $2^{\mathrm o}$ et
$3^{\mathrm o}$ énoncées ci-dessus sont bien vérifiées, mais $A$ n'est même pas
intégralement clos.

On peut varier cet exemple, et le lecteur verra sans peine que si $k$ est un corps
algébriquement clos, l'anneau $A$, localisé de l'anneau
$k[U,V,W]/(U^2-WV^2)$ en l'idéal maximal engendré par les images de $U$, $V$ et $W$,
vérifie aussi les conditions $1^{\mathrm o}$, $2^{\mathrm o}$ et $3^{\mathrm o}$
ci-dessus.

Il se pourrait que ces trois conditions impliquent même que chacun des anneaux
$B'/\mathfrak p_i'$ soit un anneau de valuation discrète, ce qui entraînerait que
$A'$ est même un anneau \emph{régulier}.
\end{env}

\begin{proposition}[21.13.10]
\label{IV.21.13.10}
Soient $X$ un préschéma localement noethérien, $Y$ une partie fermée de $X$. Pour que
le couple $(X,Y)$ soit parafactoriel il faut et il suffit que, pour tout $y\in Y$,
l'anneau local $\mathcal O_{X,y}$ soit parafactoriel.
\end{proposition}

\begin{proof}
Posons $U=X-Y$ et soit $j:U\to X$ l'injection canonique.

En vertu de (21.13.4), dire que l'homomorphisme canonique
$\mathcal O_X\to j_*(\mathcal O_U)$ soit bijectif équivaut à dire que chacun des
anneaux locaux $\mathcal O_{X,y}$ pour $y\in Y$ satisfait à la condition (i bis) de
(21.13.8).

Montrons d'abord que si le couple $(X,Y)$ est parafactoriel, chacun des anneaux
$\mathcal O_{X,y}$ ($y\in Y$) est parafactoriel. Compte tenu de la remarque
précédente, tout revient à voir que, si l'on pose
$T_y=\operatorname{Spec}(\mathcal O_{X,y})$ et $U_y=T_y-\{y\}$, tout
$\mathcal O_{U_y}$-Module inversible $\mathscr L_0$ est isomorphe à
$\mathcal O_{U_y}$. Or, lorsque $V$ parcourt l'ensemble des voisinages ouverts
affines de $y$ dans $X$, il résulte de (8.2.13) et $(\mathbf I, 2.4.2)$ que le
préschéma $U_y$ est limite projective des préschémas induits par $X$ sur les ouverts
$V-(V\cap\overline{\{y\}})$, qui sont quasi-compacts puisque $X$ est localement
noethérien. Comme les $V-(V\cap\overline{\{y\}})$ sont séparés, il résulte alors de
(8.5.2, (ii)) et (8.5.5) qu'il y a un voisinage ouvert affine $V$ de $y$ dans $X$ et
un $\mathcal O_{U_V}$-Module inversible $\mathscr L'_V$ tel que $\mathscr L_0$ soit le
faisceau induit par $\mathscr L'_V$ sur $U_y$. Mais l'hypothèse entraîne, en vertu de
(21.13.6, (i) et (iii)), que le couple $(V,V\cap\overline{\{y\}})$ est parafactoriel;
on en conclut par définition (21.13.1) qu'il existe un $\mathcal O_V$-Module inversible
$\mathscr L_V$ induisant $\mathscr L'_V$ sur $U_V$. Remplaçant au besoin $V$ par un
voisinage ouvert plus petit de $y$, on peut supposer que $\mathscr L_V$ est isomorphe
à $\mathcal O_V$, d'où on conclut que $\mathscr L_0$ est isomorphe à
$\mathcal O_{U_y}$.

Inversement, prouvons que si tous les $\mathcal O_{X,y}$ ($y\in Y$) sont
parafactoriels le couple $(X,Y)$ est parafactoriel. On est évidemment ramené, compte
tenu de la remarque du début, à montrer que pour tout $\mathcal O_U$-Module inversible
$\mathscr L$, $j_*(\mathscr L)$ est un $\mathcal O_X$-Module inversible (21.13.5).
La question étant locale sur $X$, on peut supposer $X$ noethérien. L'ensemble $V$ des
$x\in X$ tels que la restriction de $j_*(\mathscr L)$ à un voisinage ouvert de $x$
dans $X$ soit inversible est évidemment ouvert dans $X$; il s'agit de montrer que
$V=X$. Pour cela, supposons le contraire et posons $Z=X-V\neq\varnothing$. Soit $y$
un point \emph{maximal} de $Z$; puisque $Z\subset Y$ par définition,
$\mathcal O_{X,y}$ est par hypothèse un anneau parafactoriel. Remplaçant au besoin
$X$ par un voisinage ouvert de $y$ dans $X$, on peut supposer que la restriction de
$j_*(\mathscr L)$
\oldpage[IV-4]{320}
à $V-(V\cap\overline{\{y\}})$ est inversible; donc, avec les notations introduites
ci-dessus, le $\mathcal O_{U_y}$-Module $\mathscr L_0$ induit par $j_*(\mathscr L)$
sur $U_y$ est inversible. Puisque $\mathcal O_{X,y}$ est parafactoriel,
$\mathscr L_0$ est induit par un $\mathcal O_{T_y}$-Module inversible $\mathscr L'$;
mais comme $T_y$ est limite projective des voisinages ouverts $W$ de $y$ dans $X$,
on peut de nouveau appliquer (8.5.2, (ii)) et (8.5.5), prouvant l'existence d'un tel
voisinage $W$ et d'un $\mathcal O_W$-Module inversible $\mathscr L''$ induisant
$\mathscr L'$ sur $T_y$, donc induisant $\mathscr L_0$ sur $U_y$; appliquant cette
fois (8.5.2.5) et (8.5.2, (i)), on en déduit qu'en remplaçant au besoin $W$ par un
voisinage ouvert plus petit de $y$, on peut supposer que les restrictions de
$j_*(\mathscr L)$ et de $\mathscr L''$ à
$W-(W\cap\overline{\{y\}})$ soient égales. Si l'on pose alors $V_1=V\cup W$ et si
l'on désigne par $\mathscr L_1$ le $\mathcal O_{V_1}$-Module inversible égal à
$\mathscr L''$ dans $W$, à $j_*(\mathscr L)|V$ dans $V$, on a $U\subset V_1$ et
$\mathscr L_1|U=\mathscr L$. On conclut alors du fait que l'homomorphisme
$\mathcal O_X\to j_*(\mathcal O_U)$ est bijectif et de (21.13.2, $b$)) que
$j_*(\mathscr L)|V_1$ est isomorphe à $\mathscr L_1$, donc inversible. Mais cela
contredit la définition de $V$, et conclut donc la démonstration de (21.13.10).
\end{proof}

\begin{corollary}[21.13.11]
\label{IV.21.13.11}
Soient $X$ un préschéma localement noethérien, $Y$ une partie fermée de $X$,
$U=X-Y$. Supposons que le couple $(X,Y)$ soit parafactoriel et que $U$ soit
localement factoriel; en d'autres termes (21.13.10), pour tout $x\in X$, l'anneau
$\mathcal O_{X,x}$ est : $1^{\mathrm o}$ parafactoriel si $x\in Y$;
$2^{\mathrm o}$ factoriel si $x\notin Y$. Alors $X$ est localement factoriel
(autrement dit, $\mathcal O_{X,x}$ est en fait factoriel pour tout $x\in X$).
\end{corollary}

\begin{proof}
Supposons en effet le contraire, et soit $y$ un point de $Y$ tel que
$\mathcal O_{X,y}$ ne soit pas factoriel et ait une dimension \emph{minima} parmi
tous les points de $Y$ ayant cette propriété. Alors, si l'on pose
$T_y=\operatorname{Spec}(\mathcal O_{X,y})$, $U_y=T_y-\{y\}$, l'hypothèse que
$\mathcal O_{X,y}$ est parafactoriel entraîne $\operatorname{Pic}(U_y)=0$ et
$\operatorname{prof}(\mathcal O_{X,y})\geq2$ (21.13.8); de plus le choix de $y$
entraîne que $U_y$ est localement factoriel. Mais alors (21.6.14) prouve que
$\mathcal O_{X,y}$ est factoriel, contrairement à l'hypothèse.
\end{proof}

\begin{proposition}[21.13.12]
\label{IV.21.13.12}
Soient $A$, $B$ deux anneaux locaux noethériens,
$\rho:A\to B$ un homomorphisme local faisant de $B$ un $A$-module plat. Alors, si
$B$ est factoriel, il en est de même de $A$.
\end{proposition}

\begin{proof}
On sait déjà que $A$ est intègre et intégralement clos (6.5.4), donc on peut se borner
au cas où $\dim(A)\geq2$, et raisonner par récurrence sur $n=\dim(A)$. Soient
$X=\operatorname{Spec}(A)$, $a$ le point fermé de $X$,
$X'=\operatorname{Spec}(B)$, $f:X'\to X$ le morphisme structural,
$U=X-\{a\}$, $U'=f^{-1}(U)$.

Les anneaux locaux $\mathcal O_{X',x'}$ aux points $x'\in f^{-1}(a)$ sont factoriels
et de dimension $\geq2$ (6.1.2), donc parafactoriels (21.13.9, (ii)); on en conclut
(21.13.10) que le couple $(X',f^{-1}(a))$ est parafactoriel. En vertu de (21.13.6,
(iii)), cela montre que l'anneau $A$ est parafactoriel, donc
$\operatorname{prof}(A)\geq2$ et $\operatorname{Pic}(U)=0$ (21.13.8). Enfin, les
anneaux locaux de $U$ étant de dimension $<n$, l'hypothèse de récurrence prouve que
$U$ est \emph{localement factoriel}; la conclusion résulte alors de (21.6.14).
\end{proof}

\begin{env}[21.13.12.1]
\label{IV.21.13.12.1}
On notera que l'énoncé (21.13.12) où on remplace « factoriel » par « parafactoriel »
n'est plus exact, comme le montre l'exemple où $A=k$ est un corps et $B$ l'anneau
local factoriel $k[[T_1,T_2]]$. Cependant, il résulte de (21.13.6, (iii)) que, sous
les conditions de (21.13.12), si $\mathfrak m$ désigne l'idéal maximal de $A$ et si
$\mathfrak mB$ est un \emph{idéal}
\oldpage[IV-4]{321}
\emph{de définition} de $B$ (autrement dit, si $\dim(B/\mathfrak mB)=0$), alors,
lorsque $B$ est parafactoriel, il en est de même de $A$. En particulier, si le
complété $\widehat A$ d'un anneau local noethérien est parafactoriel, il en est de
même de $A$.
\end{env}

\begin{remark}[21.13.13]
\label{IV.21.13.13}
Une partie des définitions et résultats précédents se rattache à des considérations
plus générales sur la cohomologie des faisceaux de groupes commutatifs sur un espace
topologique, développées au chap.~III, 3\textsuperscript{e} partie, et nous n'en avons
donné ici un exposé indépendant que pour la commodité du lecteur. En effet, si $X$
est un espace annelé, on a une \emph{équivalence} canonique entre la catégorie des
$\mathcal O_X$-Modules inversibles et celle des \emph{faisceaux principaux homogènes
sous le faisceau de groupes} $\mathcal O_X^*$ (16.5.15). Cette équivalence se définit
en faisant correspondre à tout $\mathcal O_X$-Module inversible $\mathscr L$ le
faisceau $\mathscr L^*=\mathscr Isom_{\mathcal O_X}(\mathcal O_X,\mathscr L)$ des
germes d'isomorphismes de $\mathcal O_X$ sur $\mathscr L$; il est immédiat que
$\mathscr L^*$ est canoniquement muni d'une structure de faisceau principal homogène
sous $\mathcal O_X^*\cong\mathscr Isom_{\mathcal O_X}(\mathcal O_X,\mathcal O_X)$.
La vérification du fait que le foncteur $\mathscr L\mapsto\mathscr L^*$ est une
équivalence de catégories est immédiate.

Cela étant, considérons de façon générale un espace topologique $X$ et un faisceau de
groupes $\mathscr G$ sur $X$; soit $Y$ une partie fermée de $X$, posons $U=X-Y$, et
soit $i$ un entier tel que $0\leq i\leq2$; on dira que le couple $(X,Y)$ est
\emph{$i$-pur pour $\mathscr G$} si, pour tout ouvert $V$ de $X$, le foncteur
restriction $\mathscr P\mapsto\mathscr P|(U\cap V)$, de la catégorie de faisceaux
principaux homogènes sous le faisceau de groupes $\mathscr G|V$, dans la catégorie
analogue sous le faisceau de groupes $\mathscr G|(U\cap V)$, est \emph{fidèle}
($i=0$), resp.\ \emph{pleinement fidèle} ($i=1$), resp.\ une \emph{équivalence de
catégories} ($i=2$). Dans le cas où $X$ est un espace annelé et
$\mathscr G=\mathcal O_X^*$, dire que $(X,Y)$ est $i$-pur signifie donc pour $i=0$,
que l'homomorphisme $\mathcal O_X\to j_*(\mathcal O_U)$ est \emph{injectif}; pour
$i=1$, que cet homomorphisme est \emph{bijectif}; et enfin, pour $i=2$, que le couple
$(X,Y)$ est \emph{parafactoriel} (21.13.3).

Revenant au cas général, rappelons que les morphismes de la catégorie des faisceaux
principaux sous $\mathscr G$ sont par définition les isomorphismes. Dire que le couple
$(X,Y)$ est $0$-pur signifie donc que pour tout ouvert $V$ de $X$, l'homomorphisme
canonique $\mathrm H^0(V,\mathscr G)\to\mathrm H^0(U\cap V,\mathscr G)$ est injectif;
dire que le couple $(X,Y)$ est $1$-pur signifie que l'homomorphisme canonique
$\mathrm H^0(V,\mathscr G)\to\mathrm H^0(U\cap V,\mathscr G)$ est bijectif (et alors
l'homomorphisme canonique
$\mathrm H^1(V,\mathscr G)\to\mathrm H^1(U\cap V,\mathscr G)$ est injectif); enfin,
on peut montrer que pour que le couple $(X,Y)$ soit $2$-pur, il faut et il suffit que
les homomorphismes
$\mathrm H^i(V,\mathscr G)\to\mathrm H^i(U\cap V,\mathscr G)$ soient bijectifs pour
$i=0$ et $i=1$. Lorsque $\mathscr G$ est un faisceau de groupes commutatifs, en
introduisant les faisceaux de cohomologie $\mathscr H_Y^p(\mathscr G)$ définis au
chap.~III, 3\textsuperscript{e} partie, dire que $(X,Y)$ est $i$-pur pour $i\leq2$
signifie que $\mathscr H_Y^p(\mathscr G)=0$ pour $p\leq i$; sous cette forme, la
notion se généralise donc aussitôt pour tout entier $i$.
\end{remark}

\begin{proposition}[21.13.14]
\label{IV.21.13.14}
Soient $X$ un préschéma localement noethérien \emph{normal},
$(U_\lambda)_{\lambda\in L}$ une famille filtrante décroissante d'ouverts de $X$. Les
conditions suivantes sont équivalentes :
\begin{enumerate}[label=\emph{\alph*)}]
  \item Tout cycle $1$-codimensionnel $Z$ sur $X$ tel qu'il existe un indice
  $\lambda$ pour lequel $Z|U_\lambda$ soit localement principal (21.6.7), est
  localement principal.
  \item Pour tout $x\in X$ tel que $\dim(\mathcal O_{X,x})\geq2$ et que $x$
  n'appartienne pas à $\bigcap_{\lambda\in L}U_\lambda$, l'anneau
  $\mathcal O_{X,x}$ est parafactoriel.
\end{enumerate}
\end{proposition}

\begin{proof}
\oldpage[IV-4]{322}
La condition $b)$ peut encore s'exprimer de la façon suivante :

\emph{$b')$ Pour toute partie fermée $Y$ de $X$ telle que
$\operatorname{codim}(Y,X)\geq2$ et telle qu'il existe $\lambda$ pour lequel
$Y\subset X-U_\lambda$, le couple $(X,Y)$ est parafactoriel.}

La propriété $a)$ peut encore s'exprimer de la façon suivante : si l'ensemble fermé
$N$ des points $x\in X$ où $Z$ n'est pas principal est contenu dans l'un des
$X-U_\lambda$, alors $Z$ est localement principal. Comme $X$ est normal, la condition
$\dim(\mathcal O_{X,x})\leq1$ entraîne que $\mathcal O_{X,x}$ est un corps ou un
anneau de valuation discrète, donc $Z_x$ est principal; en d'autres termes, on a
nécessairement $\operatorname{codim}(N,X)\geq2$ (5.1.3). La propriété $a)$ est donc
équivalente à la suivante :

\emph{$a')$ S'il existe un ensemble fermé $Y$ contenu dans un des $X-U_\lambda$, tel
que $\operatorname{codim}(Y,X)\geq2$ et que $Z|(X-Y)$ soit localement principal,
alors $Z$ est localement principal.}

Notons d'autre part que $b)$ implique $b')$ en vertu de (21.13.10); inversement, si
$b')$ est vérifiée et si $x\notin U_\lambda$, alors $Y=\{x\}$ est contenue dans
$X-U_\lambda$ et on a $\operatorname{codim}(Y,X)\geq2$, donc il résulte encore de
(21.13.10) que $\mathcal O_{X,x}$ est parafactoriel, ce qui prouve l'équivalence de
$b)$ et $b')$. On est ainsi ramené à prouver l'équivalence de $a')$ et de $b')$.
Notons que puisque $X$ est normal, on a, pour toute partie fermée $Y$ de $X$ telle que
$\operatorname{codim}(Y,X)\geq2$, $\operatorname{prof}_Y(\mathcal O_X)\geq2$
(5.8.6); si l'on pose $U=X-Y$, l'homomorphisme
$\mathcal O_X\to j_*(\mathcal O_U)$ est donc bijectif (21.13.4); comme les conditions
$a')$ et $b')$ sont locales, on voit qu'il suffit de montrer l'équivalence des
conditions suivantes, lorsque $X$ est noethérien et normal et $Y$ fermé dans $X$ et
tel que $\operatorname{codim}(Y,X)\geq2$ :

\emph{$a'')$ Pour tout cycle $1$-codimensionnel $Z$ sur $X$, l'hypothèse que $Z|U$
(où $U=X-Y$) est localement principal entraîne que $Z$ lui-même est localement
principal.}

\emph{$b'')$ L'homomorphisme canonique
$\operatorname{Pic}(X)\to\operatorname{Pic}(U)$ est surjectif.}

Prouvons d'abord que $a'')$ entraîne $b'')$; un élément
$\zeta_0\in\operatorname{Pic}(U)$ a pour image dans $\mathfrak{Cl}(U)$ (21.6.11) la
classe d'un cycle $1$-codimensionnel $Z_0$ sur $U$, qui est localement principal.
L'hypothèse $\operatorname{codim}(Y,X)\geq2$ entraîne que l'homomorphisme de
restriction $\mathscr Z^1(X)\to\mathscr Z^1(U)$ est bijectif, donc $Z_0=Z|U$, où $Z$
est un cycle $1$-codimensionnel sur $X$; puisque $Z_0$ est localement principal, il en
est de même de $Z$ en vertu de $a'')$; il résulte de (21.6.11) que l'image de $Z$ dans
$\mathfrak{Cl}(X)$ est l'image d'un unique élément $\zeta\in\operatorname{Pic}(X)$,
et il est clair alors que $\zeta_0$ est l'image de $\zeta$.

Inversement, prouvons que $b'')$ entraîne $a'')$. Soit donc $Z$ un cycle
$1$-codimensionnel sur $X$ tel que $Z|U$ soit localement principal; l'image de $Z|U$
dans $\mathfrak{Cl}(U)$ est donc l'image d'un unique élément
$\zeta_0\in\operatorname{Pic}(U)$ (21.6.11). Par hypothèse il existe
$\zeta\in\operatorname{Pic}(X)$ dont l'image dans $\operatorname{Pic}(U)$ est
$\zeta_0$; l'image de $\zeta$ dans $\mathfrak{Cl}(X)$ est donc la classe d'un cycle
$1$-codimensionnel $Z'$ sur $X$ tel que $Z|U$ et $Z'|U$ soient linéairement
équivalents. Mais puisque $U$ est schématiquement dense dans $X$ l'image de
$\mathscr Z^1_{\mathrm{princ}}(X)$ par l'isomorphisme
$\mathscr Z^1(X)\to\mathscr Z^1(U)$ est $\mathscr Z^1_{\mathrm{princ}}(U)$, donc $Z$
et $Z'$ sont linéairement équivalents, et puisque $Z'$ est localement principal, il
en est de même de $Z$.
\end{proof}

\begin{corollary}[21.13.15]
\label{IV.21.13.15}
Soient $X$ un préschéma localement noethérien et normal, $S$ une partie de $X$. Les
conditions suivantes sont équivalentes :
\begin{enumerate}[label=\emph{\alph*)}]
  \item Tout cycle $1$-codimensionnel sur $X$ qui est principal aux points de $S$ est
  localement principal.
  \oldpage[IV-4]{323}
  \item Pour tout $x\in X$ qui n'est pas générisation d'un point de $S$ (autrement
  dit, tel que $\overline{\{x\}}\cap S=\varnothing$) et qui est tel que
  $\dim(\mathcal O_{X,x})\geq2$, l'anneau $\mathcal O_{X,x}$ est parafactoriel.
\end{enumerate}
\end{corollary}

\begin{proof}
Il est clair que l'ensemble $S'$ des points de $X$ générisations de points de $S$ est
l'intersection des voisinages ouverts de $S$. On peut se borner au cas où $X$ est
noethérien (les propriétés $a)$ et $b)$ étant locales sur $X$); comme tout cycle
$1$-codimensionnel sur $X$ qui est principal aux points de $S$ l'est aussi aux points
d'un voisinage ouvert de $S$, on voit que la condition $a)$ de (21.13.15) équivaut à
la condition $a)$ de (21.13.14) appliquée à la famille $(U_\lambda)$ des voisinages
ouverts de $S$. Le corollaire résulte alors de l'équivalence de $a)$ et $b)$ dans
(21.13.14).
\end{proof}

\begin{remark}[21.13.16]
\label{IV.21.13.16}
Sous les conditions générales de (21.13.14), supposons de plus que l'on ait
$\varinjlim\operatorname{Pic}(U_\lambda)=0$. Alors la condition $a)$ de (21.13.14)
est aussi équivalente à la suivante :
\begin{enumerate}
  \item[c)] \emph{Tout cycle $1$-codimensionnel sur $X$ dont le support est contenu
  dans un des ensembles $X-U_\lambda$ est localement principal.}
\end{enumerate}

On peut se borner au cas où $X$ est irréductible et tous les $U_\lambda$ sont non
vides. Il est clair que $a)$ entraîne $c)$, car si le cycle $1$-codimensionnel $Z$ sur
$X$ a son support dans $X-U_\lambda$, $Z|U_\lambda=0$, donc $Z|U_\lambda$ est
localement principal. Inversement $c)$ entraîne $a)$ : soit en effet $Z$ un cycle
$1$-codimensionnel sur $X$ tel que $Z|U_\lambda$ soit localement principal; comme $X$
est normal, $Z|U_\lambda=\operatorname{cyc}(D_\lambda)$, où $D_\lambda$ est un
diviseur sur $U_\lambda$ (21.6.10, (ii)), et l'hypothèse implique (en vertu de
(21.3.4)) qu'il y a un ensemble $U_\mu\subset U_\lambda$ tel que
$D_\mu=D_\lambda|U_\mu$ soit équivalent à $0$. Si
$D_\mu=\operatorname{div}(f_\mu)$, où $f_\mu$ est une fonction rationnelle régulière
sur $U_\mu$, $f_\mu$ peut être considérée comme une fonction rationnelle régulière
sur $X$. Si $Z'=\operatorname{cyc}(\operatorname{div}(f_\mu))$, on voit donc que
$Z''=Z-Z'$ a son support dans $X-U_\mu$; en vertu de l'hypothèse, $Z''$ est
localement principal, d'où la conclusion.

La condition $\varinjlim\operatorname{Pic}(U_\lambda)=0$ sera en particulier remplie
si $(U_\lambda)$ est la famille des voisinages ouverts d'un \emph{point} $z\in X$,
car pour tout $\mathcal O_{U_\lambda}$-Module inversible $\mathscr L_\lambda$, il
existe par définition un $U_\mu\subset U_\lambda$ tel que
$\mathscr L_\lambda|U_\mu$ soit isomorphe à $\mathcal O_{U_\mu}$. On voit donc que,
dans l'énoncé de (21.13.15), si $S=\{z\}$, les conditions $a)$ et $b)$ équivalent
aussi à la suivante :
\begin{enumerate}
  \item[c)] \emph{Tout cycle $1$-codimensionnel sur $X$ dont le support ne contient
  pas $z$ est localement principal.}
\end{enumerate}
Si de plus $\operatorname{Pic}(X)=0$, on en conclura que cette condition entraîne que
tout cycle $1$-codimensionnel dont le support ne contient pas $z$ est principal.
\end{remark}

\subsection{Le théorème de Ramanujam--Samuel}
\label{subsection:IV.21.14}
\label{IV.21.14}

\begin{theorem}[21.14.1]
\label{IV.21.14.1}
\textup{(Ramanujam--Samuel)}. ---
Soit $A$ un anneau local noethérien d'idéal maximal $\mathfrak m$, tel que son complété
$\widehat A$ soit intègre et intégralement clos. Soient $B$ un anneau local noethérien
tel que $\dim(B)\geq\dim(A)$, $\rho:A\to B$ un homomorphisme local faisant de $B$
une $A$-algèbre formellement lisse (pour les topologies préadiques
$(\mathbf 0, 19.3.1)$) et tel que le corps résiduel de $B$ soit fini sur celui de $A$.
Alors tout cycle $1$-codimensionnel sur $\operatorname{Spec}(B)$ qui est principal au
point $\mathfrak p=\mathfrak mB$, est un cycle $1$-codimensionnel principal.
\end{theorem}

\begin{proof}
Si $k=A/\mathfrak m$ est le corps résiduel de $A$, $B/\mathfrak mB$ est une
$k$-algèbre formellement lisse (pour sa topologie préadique) $(\mathbf 0, 19.3.5)$,
\oldpage[IV-4]{324}
donc régulière, et en particulier intègre, autrement dit
$\mathfrak p=\mathfrak mB$ est bien un idéal premier de $B$, ce qui justifie
l'énoncé. Tout revient évidemment à prouver que tout idéal premier $\mathfrak q$ de
$B$ non contenu dans $\mathfrak p$ est \emph{principal}.

Soient $\widehat A$, $\widehat B$ les complétés de $A$ et $B$ respectivement, de
sorte que l'idéal maximal de $\widehat A$ est $\mathfrak m\widehat A$; on sait
$(\mathbf 0, 19.3.6)$ que $\widehat B$ est une $\widehat A$-algèbre formellement
lisse pour les topologies adiques. Soit $k'$ le corps résiduel de $\widehat B$,
extension finie de $k$; il existe un homomorphisme local $\widehat A\to C$, où $C$
est un anneau local noethérien qui est un $\widehat A$-module \emph{fini} et
\emph{plat} (donc libre) et est tel que $C/\mathfrak mC$ soit isomorphe à $k'$
$(\mathbf 0_{\mathrm{III}}, 10.3.1)$; on en déduit que $C$ est \emph{complet}, et il
résulte alors de (7.5.1), (7.5.3) et (6.5.4, (ii)) que $C$ est \emph{intègre et
intégralement clos}. En outre,
\[
  D=\widehat B\widehat\otimes_{\widehat A}C
\]
est un anneau semi-local complet, composé direct d'anneaux locaux complets dont l'un,
$D_0$, a pour corps résiduel $k'$ (puisque $k'\otimes_k k'$ est composé direct
d'anneaux locaux dont l'un est isomorphe à $k'$). Comme $D$ est formellement lisse
sur $C$, il en est de même de $D_0$; par suite $D_0/\mathfrak mD_0$ est une
$k'$-algèbre formellement lisse, de corps résiduel $k'$, ce qui entraîne qu'elle est
$k'$-isomorphe à une algèbre de séries formelles
$k'[[T_1,\ldots,T_n]]$ $(\mathbf 0, 19.6.4)$; on en conclut, par
$(\mathbf 0, 19.7.1.5)$, que $D_0$ est $C$-isomorphe à
$C[[T_1,\ldots,T_n]]$, et par suite \emph{intègre et intégralement clos} (Bourbaki,
\emph{Alg. comm.}, chap.~V, \S1, n\textsuperscript{o}~4, prop.~14). Comme les
morphismes $\operatorname{Spec}(D_0)\to\operatorname{Spec}(\widehat B)$ et
$\operatorname{Spec}(D_0)\to\operatorname{Spec}(B)$ sont fidèlement plats, on en
déduit que $B$ et $\widehat B$ sont aussi \emph{intègres et intégralement clos}
(2.1.13). Cela prouve que l'idéal $\mathfrak qD_0$ est divisoriel (Bourbaki,
\emph{Alg. comm.}, chap.~VII, \S1, n\textsuperscript{o}~10, prop.~15) et non contenu
dans $\mathfrak mD_0$, sans quoi on aurait
\[
  \mathfrak q=(\mathfrak qD_0)\cap B
  \subset(\mathfrak mD_0)\cap B=\mathfrak mB=\mathfrak p
\]
par fidèle platitude $(\mathbf 0_{\mathrm I}, 6.5.1)$. On en conclut qu'aucun des
idéaux premiers $\mathfrak r_h$ de hauteur $1$ dans $D_0$ qui contiennent
$\mathfrak qD_0$ ne peut être contenu dans $\mathfrak mD_0$; si l'on prouve que ces
idéaux sont principaux, il en résultera que $\mathfrak qD_0$ est principal, les
diviseurs (au sens de Bourbaki) de $\mathfrak qD_0$ et d'un produit de puissances des
$\mathfrak r_h$ étant égaux (Bourbaki, \emph{Alg. comm.}, chap.~VII, \S1,
n\textsuperscript{o}~4, prop.~5). Comme $\mathfrak qD_0\cap B=\mathfrak q$, on en
déduit par fidèle platitude que $\mathfrak q$ est principal (2.5.2), en utilisant le
fait que $B$ est un anneau local.

On est ainsi ramené à prouver le théorème dans le cas particulier où $A$ est
\emph{complet} et $B=A[[T_1,\ldots,T_n]]$. Montrons d'abord qu'on peut se ramener au
cas où $n=1$. En effet, procédons par récurrence sur $n$, et soit
$f\in\mathfrak q$ un élément n'appartenant pas à $\mathfrak p=\mathfrak mB$; il
suffira de montrer qu'il existe un $A$-automorphisme $\sigma$ de $B$ tel que
$\sigma(f)$ n'appartienne pas à l'idéal
\[
  \mathfrak p'=\mathfrak p+BT_1+\cdots+BT_{n-1};
\]
en effet, si $B'=A[[T_1,\ldots,T_{n-1}]]$, $B'$ est complet et a pour idéal maximal
$\mathfrak mB'+B'T_1+\cdots+B'T_{n-1}=\mathfrak m'$. Comme on a
$B=B'[[T_n]]$ et $\mathfrak p'=\mathfrak m'B$, il résultera de l'hypothèse de
récurrence (compte tenu de ce que $B'$ est intègre et intégralement clos) que l'idéal
$\sigma(\mathfrak q)$ est principal dans $B$, donc aussi $\mathfrak q$. Or, si
$\overline f\in k[[T_1,\ldots,T_n]]$ est l'image de $f$ (« série réduite » de $f$),
on sait (Bourbaki, \emph{Alg. comm.}, chap.~VII, \S3, n\textsuperscript{o}~7,
lemme~3) qu'il y a un $A$-automorphisme $\sigma$ de $B$ tel que
$(\overline{\sigma(f)})(0,\ldots,0,T_n)\neq0$, ce qui implique évidemment
$\sigma(f)\notin\mathfrak p'$.

Supposons donc $n=1$ et écrivons $T$ au lieu de $T_1$, de sorte que $B=A[[T]]$. Il
suffit de montrer que dans l'anneau de polynômes $A[T]$, l'idéal
$\mathfrak q\cap A[T]$ est principal;
\oldpage[IV-4]{325}
en effet, soit $f_0$ un élément de $\mathfrak q$ n'appartenant pas à $\mathfrak mB$;
c'est une série formelle dont la série réduite $\overline f_0$ n'est pas nulle; donc,
en vertu du théorème de préparation (Bourbaki, \emph{Alg. comm.}, chap.~VII, \S3,
n\textsuperscript{o}~8, prop.~5), pour tout $f\in\mathfrak q$, il existe $g\in B$ et
un polynôme $r\in A[T]$ tels que $f=gf_0+r$ et on a donc
$r\in\mathfrak q\cap A[T]$; d'autre part (\emph{loc. cit.}, prop.~6) il existe un
polynôme distingué non constant $F_0\in A[T]$ et un élément inversible $u\in B$ tels
que $f_0=uF_0$, donc on a aussi $F_0\in\mathfrak q\cap A[T]$, ce qui prouve que
$\mathfrak q$ est engendré par $\mathfrak q\cap A[T]=\mathfrak q_1$. Puisque $B$
est plat sur $A[T]$ $(\mathbf 0_{\mathrm I}, 7.3.3)$, il résulte de Bourbaki,
\emph{Alg. comm.}, chap.~VII, \S1, n\textsuperscript{o}~10, prop.~15, que
$\mathfrak q_1$ est un idéal premier de hauteur $1$ dans $A[T]$.

En outre, on a nécessairement $\mathfrak q_1\cap A=0$; sinon,
$\mathfrak q_1\cap A$ serait nécessairement de hauteur $\geq1$, et il résulterait de
(5.5.3) que l'on aurait $\mathfrak q_1=(\mathfrak q_1\cap A)A[T]$. Mais alors,
puisque $\mathfrak q_1\cap A\subset\mathfrak m$, on aurait
$\mathfrak q_1\subset\mathfrak mA[T]$ contrairement à l'hypothèse sur
$\mathfrak q$. Si $K$ est le corps des fractions de $A$, $\mathfrak q_1K[T]$ est
donc un idéal premier distinct de $0$ et de $K[T]$ dans $K[T]$, donc de la forme
$hK[T]$, où
\[
  h(T)=T^m+a_1T^{m-1}+\cdots+a_m,
  \qquad m\geq1,
\]
avec les $a_i\in K$, $h$ étant irréductible dans $K[T]$. Mais on a vu plus haut qu'il
existe dans $\mathfrak q_1$ un polynôme distingué non constant $F_0\in A[T]$. Si $t$
est la classe de $T$ dans $K[T]/\mathfrak q_1K[T]$, $t$ est donc racine du polynôme
$F_0$ dans une extension de $K$ et par suite $h$ divise $F_0$ dans $K[T]$; mais
puisque $h$ et $F_0$ sont unitaires, cela entraîne que les coefficients $a_i$ de $h$
sont entiers sur $A$ (Bourbaki, \emph{Alg. comm.}, chap.~V, \S1,
n\textsuperscript{o}~3, prop.~11), donc appartiennent à $A$ puisque $A$ est
intégralement clos. Autrement dit, on a
\[
  h\in A[T]\cap\mathfrak q_1K[T]=\mathfrak q_1
\]
(Bourbaki, \emph{Alg. comm.}, chap.~II, \S1, n\textsuperscript{o}~5, prop.~11);
comme tout polynôme $g\in\mathfrak q_1$ est divisible par $h$ dans $K[T]$ et que $h$
est unitaire, les coefficients de $g/h$ appartiennent à $A$, donc
$\mathfrak q_1=hA[T]$. C.Q.F.D.
\end{proof}

L'énoncé (21.14.1) est \emph{équivalent} au suivant :

\begin{corollary}[21.14.2]
\label{IV.21.14.2}
Sous les hypothèses de (21.14.1) sur $A$, $B$ et $\rho$ pour tout idéal premier
$\mathfrak q$ de $B$ non contenu dans $\mathfrak p=\mathfrak mB$ et tel que
$\dim(B_\mathfrak q)\geq2$, l'anneau $B_\mathfrak q$ est parafactoriel. En
particulier, si $\dim(B\otimes_A k)>0$ (\emph{i.e.} $\dim B>\dim A$
$((\mathbf 0, 19.7.1)$ et $(\mathbf 0, 6.1.1)))$, l'anneau $B$ est parafactoriel.
\end{corollary}

\begin{proof}
En effet, on a vu dans la démonstration de (21.14.1) que les conditions de l'énoncé
entraînent que $B$ est intègre et intégralement clos; l'équivalence des énoncés
(21.14.1) et (21.14.2) résulte alors de (21.13.15) appliqué à
$X=\operatorname{Spec}(B)$ et $S=\{\mathfrak p\}$.
\end{proof}

\begin{proposition}[21.14.3]
\label{IV.21.14.3}
Soient $S$ un préschéma normal, $f:X\to S$ un morphisme lisse.
\begin{enumerate}[label=(\roman*)]
  \item Si $S$ est localement noethérien (donc aussi $X$), alors, pour tout $x\in X$
  tel que $\dim(\mathcal O_{X,x})\geq2$ et que $x$ ne soit pas point maximal de sa
  fibre $f^{-1}(f(x))$, l'anneau $\mathcal O_{X,x}$ est parafactoriel. Tout cycle
  $1$-codimensionnel $Z$ sur $X$ tel que $\operatorname{Supp}(Z)$ ne contienne aucune
  composante irréductible d'une fibre $f^{-1}(s)=X_s$, est localement principal.
  \item Soit $Y$ une partie fermée de $X$ ne contenant aucune composante irréductible
  d'une fibre $X_s$, et telle que, pour tout point maximal $\eta$ de $S$,
  $Y_\eta=Y\cap X_\eta$ soit de codimension $\geq2$ dans $X_\eta$. Alors le couple
  $(X,Y)$ est parafactoriel.
\end{enumerate}
\end{proposition}

\begin{proof}
On notera que les conditions de (ii) sont remplies si $Y$ est une partie fermée de $X$
telle que, pour tout $s\in S$, $Y_s=Y\cap X_s$ soit de codimension $\geq2$ dans
$X_s$.

Pour prouver (21.14.3), notons d'abord que sous les hypothèses de (i), si l'on pose
$Y=\overline{\{x\}}$, il existe un voisinage ouvert $V$ de $x$ dans $X$ tel que
$Y\cap V$ et $V$ vérifient les conditions de (ii) : en effet, il résulte de l'hypothèse
et de (9.5.3) que l'on peut prendre $V$ tel que $V\cap Y$ ne contienne aucune
composante irréductible d'un $X_s$; d'autre part, pour tout point maximal $\eta$ de
$S$ tel que $Y\cap X_\eta\neq\varnothing$, les points de $Y\cap X_\eta$ sont des
spécialisations de $x$, donc en un tel point $z$ on a
$\dim(\mathcal O_{X,z})\geq2$, et comme
$\mathcal O_{X,z}=\mathcal O_{X_\eta,z}$ ($S$ étant réduit au point $\eta$ en vertu
de (2.1.13)), $Y\cap X_\eta$ est bien de codimension $\geq2$ dans $X_\eta$. Comme le
couple $(V,Y\cap V)$ est alors parafactoriel en vertu de (ii), on conclut de (21.13.10)
que l'anneau $\mathcal O_{X,x}$ est parafactoriel, c'est-à-dire la première assertion
de (i). Pour prouver la seconde, on peut se borner au cas où
$Z=\overline{\{z\}}$, où $z\in X^{(1)}$; comme $\mathcal O_{X,z}$ est un anneau local
intégralement clos et de dimension $1$, c'est un anneau de valuation discrète et $Z$
est donc principal au point $z$; pour tout autre point $x$ de
$\operatorname{Supp}(Z)$, on a donc $\dim(\mathcal O_{X,x})\geq2$ et $x$ n'est pas
point maximal de sa fibre par hypothèse, donc $\mathcal O_{X,x}$ est parafactoriel.
Appliquons alors (21.13.15) en prenant pour $S$ l'ensemble formé de $z$ et des points
maximaux des fibres de $f$; il est clair que $Z$ est principal aux points de $S$
puisque $z$ est le seul point de $S$ appartenant à $\operatorname{Supp}(Z)$; comme
d'autre part la condition $b)$ de (21.13.15) est évidemment remplie, on en conclut que
$Z$ est localement principal.

On est donc ramené à prouver (ii). Posons $U=X-Y$ et soit $j:U\to X$ l'injection
canonique.
\oldpage[IV-4]{326}
Comme les hypothèses et la conclusion sont locales sur $S$ et sur $X$ en vertu de
(21.13.6, (i)), on peut se borner au cas où $S$ et $X$ sont affines, $f$ étant donc de
présentation finie. Puisque les fibres de $f$ sont régulières, l'hypothèse sur $Y$
entraîne que, pour tout point $x\in Y$, on a
\[
  \operatorname{prof}(\mathcal O_{X_{f(x)},x})
  =\dim(\mathcal O_{X_{f(x)},x})\geq1;
\]
il résulte de (19.9.8) que l'homomorphisme canonique
$\mathcal O_X\to j_*(\mathcal O_U)$ est \emph{injectif}. De plus, remplaçant $Y$ par
une partie fermée plus grande suivant la méthode de la démonstration de (19.9.8), et
utilisant (21.13.6, (ii)), on peut supposer que $Y$ est défini par un idéal \emph{de
type fini} de l'anneau de $X$, donc l'ouvert $U=X-Y$ est \emph{quasi-compact et
quasi-séparé} et l'ensemble fermé $Y$ est \emph{constructible}.

Pour prouver (ii), il suffit, pour tout $\mathcal O_U$-Module inversible $\mathscr L$
et tout point $x\in Y$ donnés, d'établir l'existence d'un voisinage ouvert $V$ de $x$
dans $X$ tel que : $1^{\mathrm o}$ l'homomorphisme canonique
$\mathcal O_V\to(j_V)_*(\mathcal O_{U\cap V})$ est surjectif;
$2^{\mathrm o}$ il existe un $\mathcal O_V$-Module inversible $\mathscr L'_V$ tel que
$\mathscr L'_V|(U\cap V)=\mathscr L|(U\cap V)$ ((21.13.5) et (21.13.3)).

Posons $s=f(x)$; nous allons voir qu'on peut se borner au cas où
$S=S_0=\operatorname{Spec}(\mathcal O_{S,s})$. Soit en effet $(S_\nu)$ un système
fondamental de voisinages ouverts affines de $s$ dans $S$, et posons
$X_\nu=f^{-1}(S_\nu)$, $Y_\nu=Y\cap X_\nu$, $U_\nu=X_\nu-Y_\nu$,
$j_\nu:U_\nu\to X_\nu$ étant l'injection canonique; $X_0=X\times_S S_0$ est alors
limite projective des $X_\nu$ (8.1.2, $a$)); supposons la proposition vraie pour le
morphisme $f_0:X_0\to S_0$ et pour $Y_0=Y\cap X_0$, et posons $U_0=X_0-Y_0$,
$j_0:U_0\to X_0$ étant l'injection canonique. La projection
$p_\nu:X_0\to X_\nu$ étant un morphisme affine, on a
\[
  (j_0)_*(\mathcal O_{U_0})
  =p_\nu^*((j_\nu)_*(\mathcal O_{U_\nu}))
  \qquad(\mathbf{II}, 1.5.2)
\]
et l'homomorphisme canonique
$\rho_0:\mathcal O_{X_0}\to(j_0)_*(\mathcal O_{U_0})$ n'est autre que
$p_\nu^*(\rho_\nu)$, où
$\rho_\nu:\mathcal O_{X_\nu}\to(j_\nu)_*(\mathcal O_{U_\nu})$ est l'homomorphisme
canonique; comme par hypothèse $\rho_0$ est surjectif, il en est de même de $\rho_\nu$
pour $\nu$ assez grand (8.5.7). Soit d'autre part $\mathscr L_0$ le
$\mathcal O_{U_0}$-Module inversible restriction de $\mathscr L$, et notons que les
$X_\nu$ sont affines, donc quasi-compacts et quasi-séparés; puisque par hypothèse il
existe un $\mathcal O_{X_0}$-Module inversible $\mathscr L'_0$ tel que
$\mathscr L'_0|U_0=\mathscr L_0$, il existe un $\nu$ assez grand et un
$\mathcal O_{X_\nu}$-Module quasi-cohérent $\mathscr L'_\nu$ de présentation finie
tels que $\mathscr L'_0$ soit isomorphe à $p_\nu^*(\mathscr L'_\nu)$ (8.5.2, (ii));
de plus (8.5.5) on peut supposer que $\mathscr L'_\nu$ est inversible. Enfin, comme
les $U_\nu$ sont quasi-compacts et quasi-séparés et que
$\mathscr L'_0|U_0=\mathscr L_0$, il existe $\nu$ assez grand tel que
$\mathscr L'_\nu|U_\nu$ et $\mathscr L|U_\nu$ soient isomorphes (8.5.2.5).

\oldpage[IV-4]{327}
On est ainsi ramené au cas où $S=\operatorname{Spec}(A)$,
$X=\operatorname{Spec}(B)$, où $A$ est un anneau \emph{local}; puisque $X$ est
normal et $f$ lisse, $S$ est normal (17.5.7), donc $A$ est \emph{intègre et
intégralement clos}. Considérons alors $A$ comme limite inductive de ses
\emph{sous-$\mathbb Z$-algèbres de type fini}; comme la clôture intégrale d'une telle
sous-$\mathbb Z$-algèbre est un sous-anneau de $A$ et est aussi une
$\mathbb Z$-algèbre de type fini (7.8.3, (ii), (iii) et (vi)), $A$ est limite inductive
filtrante des sous-$\mathbb Z$-algèbres de type fini $A_\lambda$ de $A$ qui sont des
anneaux \emph{intégralement clos}. En vertu de (1.8.4.2), il existe un indice
$\lambda$ et une $A_\lambda$-algèbre de type fini $B_\lambda$ tels que
\[
  B=B_\lambda\otimes_{A_\lambda}A
\]
à un $A$-isomorphisme près. Posons $S_\lambda=\operatorname{Spec}(A_\lambda)$,
$X_\lambda=\operatorname{Spec}(B_\lambda)$; si $p_\lambda:X\to X_\lambda$ est la
projection canonique, on peut en outre supposer (puisque $Y$ est
\emph{constructible}) que $Y=p_\lambda^{-1}(Y_\lambda)$, où $Y_\lambda$ est une
partie fermée de $X_\lambda$ (8.3.11). Soit $f_\lambda$ le morphisme
$X_\lambda\to S_\lambda$; en vertu de la transitivité des fibres et de (4.2.6),
$Y_\lambda$ ne contient aucune composante irréductible des fibres
$f_\lambda^{-1}(s_\lambda)$ pour aucun $s_\lambda\in S_\lambda$. Puisque $f$ est
lisse, on peut aussi supposer que $f_\lambda$ est lisse (17.7.8); enfin, l'image du
point générique $\eta$ de $S$ dans $S_\lambda$ est le point générique $\eta_\lambda$
de $S_\lambda$, et on a
\[
  \operatorname{codim}((Y_\lambda)_{\eta_\lambda},
  (X_\lambda)_{\eta_\lambda})
  =\operatorname{codim}(Y_\eta,X_\eta)\geq2
\]
en vertu de (6.1.4). On voit donc que $S_\lambda$, $X_\lambda$ et $Y_\lambda$
vérifient toutes les hypothèses de (ii), et en posant
$X_\mu=X_\lambda\times_{S_\lambda}S_\mu$ pour $\lambda\leq\mu$, il en est de même de
$S_\mu$, $X_\mu$ et $Y_\mu$. Montrons que si l'on prouve que le couple
$(X_\mu,Y_\mu)$ est parafactoriel pour tout $\mu\gg\lambda$, il en est de même de
$(X,Y)$. En effet, soient $U=X-Y$, $U_\mu=X_\mu-Y_\mu$, $j:U\to X$,
$j_\mu:U_\mu\to X_\mu$ les injections canoniques; la projection
$p_\mu:X\to X_\mu$ étant un morphisme affine, on a
\[
  j_*(\mathcal O_U)=p_\mu^*((j_\mu)_*(\mathcal O_{U_\mu}))
  \qquad(\mathbf{II}, 1.5.2),
\]
et par suite l'homomorphisme canonique
$\rho:\mathcal O_X\to j_*(\mathcal O_U)$ n'est autre que $p_\mu^*(\rho_\mu)$, où
$\rho_\mu:\mathcal O_{X_\mu}\to(j_\mu)_*(\mathcal O_{U_\mu})$ est l'homomorphisme
canonique; ce dernier étant par hypothèse bijectif, il en est de même de $\rho$.
D'autre part, pour tout $\mathcal O_U$-Module inversible $\mathscr L$, il existe un
$\mu\gg\lambda$ et un $\mathcal O_{U_\mu}$-Module inversible $\mathscr L_\mu$ tels
que $\mathscr L=p_\mu'^*(\mathscr L_\mu)$ (en désignant par
$p_\mu':U\to U_\mu$ la restriction de $p_\mu$) (8.5.2 et 8.5.5), $U_\lambda$ étant
quasi-compact puisque $X_\lambda$ est noethérien. Par hypothèse, il existe un
$\mathcal O_{X_\mu}$-Module inversible $\mathscr L'_\mu$ tel que
$\mathscr L'_\mu|U_\mu$ soit isomorphe à $\mathscr L_\mu$;
$\mathscr L'=p_\mu^*(\mathscr L'_\mu)$ est alors un $\mathcal O_X$-Module
inversible tel que $\mathscr L'|U$ soit isomorphe à $\mathscr L$, ce qui prouve notre
assertion.

On est ainsi ramené à prouver (ii) lorsque l'anneau $A$ est une
$\mathbb Z$-algèbre de type fini intégralement close; comme les anneaux locaux
$\mathcal O_{S,s}$ de $S$ sont alors des anneaux excellents intégralement clos, leurs
\emph{complétés} sont aussi intégralement clos (7.8.3, (ii), (iii) et (vii)). Pour
prouver que le couple $(X,Y)$ est parafactoriel, il suffit de montrer que pour tout
$x\in Y$, l'anneau $\mathcal O_{X,x}$ est parafactoriel (21.13.10). Soit $y$ un point
fermé de $Y_{f(x)}$, qui soit spécialisation de $x$ (5.1.11); si l'on pose
$s=f(x)=f(y)$, $\mathcal O_{S,s}$ a un
\oldpage[IV-4]{328}
complété intégralement clos, et $\mathcal O_{X,y}$ est une
$\mathcal O_{S,s}$-algèbre formellement lisse pour les topologies préadiques
(17.5.3); puisque $Y$ ne contient aucune composante irréductible de $X_s$, on a
$\dim(\mathcal O_{X,y})>\dim(\mathcal O_{S,s})$ (17.5.8). Si $s\neq\eta$, on a
$\dim(\mathcal O_{S,s})\geq1$; si au contraire $s=\eta$, on a par hypothèse
$\dim(\mathcal O_{X,x})\geq2$; donc dans tous les cas
$\dim(\mathcal O_{X,y})\geq2$. En outre, l'idéal premier de $\mathcal O_{X,y}$
correspondant au point $x$ n'est pas contenu dans
$\mathfrak m_s\mathcal O_{X,y}$, puisque $Y$ ne contient aucune composante
irréductible de $X_s$. Enfin, puisque $y$ est fermé dans $X_s$, $k(y)$ est extension
finie de $k(s)$ $(\mathbf I, 6.4.2)$. Dans tous les cas, on peut appliquer à
$A=\mathcal O_{S,s}$, $B=\mathcal O_{X,y}$ et
$B_\mathfrak q=\mathcal O_{X,x}$ le résultat de (21.14.2), ce qui achève la
démonstration.
\end{proof}

\begin{env}[21.14.4]
\label{IV.21.14.4}
\emph{Remarques}. ---

(i) Il se peut que les énoncés (21.14.1) et (21.14.2) restent valables \emph{sans
hypothèse} sur le corps résiduel de $B$. Énoncé avec cette généralité, le résultat
équivaudrait, en vertu de (21.13.15) et (21.13.12.1) au suivant : Soient $A$ un
anneau local noethérien complet, intègre et intégralement clos, $B$ un anneau local
noethérien qui est une $A$-algèbre formellement lisse, tel que
$\dim(B)>\dim(A)$; alors $B$ est parafactoriel.

(ii) Nous verrons au chap.~VI, en employant une technique de « descente finie »
appliquée au morphisme $Y'\to Y=\operatorname{Spec}(A)$, où $Y'$ est le normalisé de
$Y$, que la conclusion de (21.14.1) (ou de (21.14.2)) reste valable en remplaçant
l'hypothèse que $\widehat A$ est intégralement clos par l'hypothèse que
$\widehat A$ est \emph{réduit}, pourvu que $\dim B\geq\dim A+2$. Si on remplace cette
dernière condition par $\dim B\geq\dim A+3$, on peut même supprimer l'hypothèse que
$\widehat A$ est réduit. De même, la conclusion de (21.14.3) reste valable en
remplaçant l'hypothèse que $X$ est normal par l'hypothèse que $X$ est réduit, pourvu
que $\dim(\mathcal O_{X,x})\geq\dim(\mathcal O_{S,f(x)})+2$.

(iii) Au chap.~III, 3\textsuperscript{e} partie, on prouve que si $f:X\to S$ est un
morphisme lisse, $Y$ une partie fermée localement constructible de $X$ telle que,
pour tout $s\in S$, on ait (avec les notations de (21.14.3))
$\operatorname{codim}(Y_s,X_s)\geq3$, alors le couple $(X,Y)$ est parafactoriel
([41], XII, 4.8). Cette conclusion n'est plus valable lorsqu'on suppose seulement que
$\operatorname{codim}(Y_s,X_s)\geq2$ pour tout $s\in S$ et lorsqu'on ne suppose plus
$X$ réduit. Par exemple, soient $k$ un corps, $A=k[T]/(T^2)$, algèbre des nombres
duaux sur $k$, $S=\operatorname{Spec}(A)$,
$X=\operatorname{Spec}(A[T_1,T_2])$ ($T_1$, $T_2$ indéterminées), de sorte que
$X=\mathbb V_S^2$, $Y$ étant la « section nulle » de ce fibré; si $s$ est l'unique
point fermé de $S$, on a $X_s=\mathbb V_k^2$ et $Y_s$ est la « section nulle » de
$X_s$. Pour voir que le couple $(X,Y)$ n'est pas parafactoriel, il suffit de montrer
que l'anneau $B=A[T_1,T_2]_\mathfrak m$, où $\mathfrak m$ est l'idéal
$(T_1)+(T_2)$ qui définit $Y$, n'est pas parafactoriel (21.13.10). Soit
$B_0=B_{\mathrm{red}}$, et notons $U$ et $U_0$ les complémentaires du point fermé
dans $\operatorname{Spec}(B)$ et $\operatorname{Spec}(B_0)$; en raisonnant comme
dans (21.13.9), on a la suite exacte, extension de (21.13.9.1)
\[
  \Gamma(U,\mathcal O_U)^*\longrightarrow\Gamma(U_0,\mathcal O_{U_0})^*
  \longrightarrow\mathrm H^1(U_0,\mathcal O_{U_0})
  \longrightarrow\operatorname{Pic}(U)\longrightarrow\operatorname{Pic}(U_0).
\]
Or, on a $\Gamma(U,\mathcal O_U)=B$, $\Gamma(U_0,\mathcal O_{U_0})=B_0$ puisque
$\operatorname{prof}(B)=\operatorname{prof}(B_0)=2$ (5.10.5). Par ailleurs
$\operatorname{Pic}(U_0)=0$ puisque $B_0$ est factoriel, et
$\mathrm H^1(U_0,\mathcal O_{U_0})\neq0$ [41, 3, Exemple III-1]; comme
l'homomorphisme $B^*\to B_0^*$ est surjectif, on en conclut que
$\operatorname{Pic}(U)\neq0$, donc que $B$ n'est pas parafactoriel.
\oldpage[IV-4]{329}
(iv) Le résultat (21.14.3) a d'abord été démontré par C.~Seshadri [44] dans le cas
particulier où $S$ est un schéma algébrique normal sur un corps algébriquement clos
$k$ et $X=S\times_k T$, où $T$ est un $k$-préschéma algébrique lisse sur $k$. La
démonstration de Seshadri [44, p.~188--189] est de nature globale et utilise la
théorie des schémas de Picard. Elle donne en outre (\emph{loc. cit.}) des résultats
tels que le suivant (pour lequel on ne possède pas à l'heure actuelle de démonstration
par voie locale). Soient $S$, $T$ deux préschémas localement de type fini sur un
corps $k$, $X=S\times_k T$, $Z$ un cycle $1$-codimensionnel sur $X$ (considéré comme
$S$-préschéma); supposons vérifiées les conditions suivantes :
\begin{enumerate}[label=\arabic*\textsuperscript{o}]
  \item $S$ et $T$ sont géométriquement normaux sur $k$ (6.7.6);
  \item Pour tout point maximal $\eta$ de $S$, le cycle $1$-codimensionnel $Z_\eta$
  sur la fibre $X_\eta$, ayant même multiplicité que $Z$ en tout point de
  $X^{(1)}\cap X_\eta$, est localement principal (autrement dit, est l'image d'un
  diviseur de $X_\eta$, puisque $X_\eta$ est normal);
  \item Pour tout $s\in S$, $Z$ est principal aux points maximaux de la fibre $X_s$.
\end{enumerate}
Alors $Z$ est localement principal. En d'autres termes, $X$ étant normal (6.14.1),
pour tout $x\in X$ qui n'est pas maximal dans sa fibre $X_s$ et qui n'appartient à
aucune des « fibres génériques » $X_\eta$ (ce qui implique
$\dim(\mathcal O_{X,x})\geq2$ en vertu de (6.1.1)), l'anneau local
$\mathcal O_{X,x}$ est parafactoriel, en vertu de (21.12.15).
\footnote{En fait, dans l'article précité, Seshadri suppose que $k$ est algébriquement
clos, $T$ séparé et « semi-complet » (i.e.\ tel que
$\Gamma(T,\mathcal O_T)$ soit $k$-isomorphe à $k$) et remplace l'hypothèse
$3^{\mathrm o}$ par l'hypothèse plus forte que $\operatorname{Supp}(Z)$ ne contient
aucune des fibres $X_s$ pour $s\in S$. Mais comme l'énoncé est local sur $S$, on en
conclut aussitôt qu'il suffit de faire l'hypothèse $3^{\mathrm o}$, et cela prouve que
la conclusion (interprétée comme ci-dessus en termes de propriété de parafactorialité
des anneaux $\mathcal O_{X,x}$) est locale sur $S$ et sur $T$, ce qui permet
d'éliminer complètement l'hypothèse que $T$ est « semi-complet » et que $k$ est
algébriquement clos, puisque (quitte à passer d'abord à la clôture algébrique de $k$)
on peut supposer d'abord $T$ affine, ce qui permet de le plonger comme ouvert d'un
schéma projectif et normal sur $k$, auquel le résultat de Seshadri s'applique. Notons
aussi que, grâce à cette réduction, il suffit de faire la démonstration de Seshadri dans
le cas où $T$ est \emph{projectif} (et non seulement « semi-complet »), cas où la
théorie du schéma de Picard utilisée par Seshadri est contenue dans celle qui sera
développée au chap.~VI de notre Traité.}
\end{env}

\subsection{Diviseurs relatifs}
\label{subsection:IV.21.15}
\label{IV.21.15}

\begin{env}[21.15.1]
\label{IV.21.15.1}
Soient $S$ un préschéma, $f:X\to S$ un morphisme \emph{plat et localement de
présentation finie}. On a défini dans (20.6.1) le faisceau d'anneaux
$\mathscr M_{X/S}$ des \emph{germes de fonctions méromorphes sur $X$ relatives à
$S$}, sous-faisceau de $\mathscr M_X$; il est clair que l'injection canonique
$\mathcal O_X\to\mathscr M_X$ (20.1.4.1) applique $\mathcal O_X$ sur un
sous-faisceau de $\mathscr M_{X/S}$, avec lequel on l'identifie. Soit
$\mathscr M_{X/S}^*$ le faisceau (en groupes multiplicatifs) des germes de sections
inversibles de $\mathscr M_{X/S}$ qui est donc un sous-faisceau de $\mathscr M_X^*$
et contient $\mathcal O_X^*$ comme sous-faisceau.
\end{env}

\begin{definition}[21.15.2]
\label{IV.21.15.2}
On appelle \emph{faisceau des diviseurs sur $X$ relativement à $S$}, ou
\emph{faisceau des diviseurs sur $X$ transversaux à $f$}, et on note
$\mathscr{D}iv_{X/S}$ le faisceau quotient (de groupes commutatifs)
$\mathscr M_{X/S}^*/\mathcal O_X^*$; les sections de ce faisceau au-dessus de $X$ se
nomment \emph{diviseurs sur $X$ relatifs à $S$}, ou \emph{diviseurs sur $X$
transversaux à $f$}; ils forment un groupe commutatif noté $\operatorname{Div}(X/S)$.
\end{definition}

Il est clair que $\mathscr{D}iv_{X/S}$ s'identifie canoniquement à un sous-faisceau de
$\mathscr{D}iv_X$, et par suite $\operatorname{Div}(X/S)$ à un sous-groupe de
$\operatorname{Div}(X)$, que l'on note encore \emph{additivement}. Pour
\oldpage[IV-4]{330}
tout ouvert $U$ de $X$, on a
$\mathscr M_{X/S}^*|U=\mathscr M_{U/S}^*$, donc
$\mathscr{D}iv_{X/S}|U=\mathscr{D}iv_{U/S}$, et le faisceau
$\mathscr{D}iv_{X/S}$ est donc égal au préfaisceau
$U\mapsto\operatorname{Div}(U/S)$.

Puisque $\mathscr M_{X/S}^*$ est un sous-faisceau de $\mathscr M_X^*$, les
définitions, notations et formules relatives aux diviseurs des sections de
$\mathscr M_X^*$ au-dessus de $X$ (21.1.3) s'appliquent sans changement aux sections
de $\mathscr M_{X/S}^*$ au-dessus de $X$.

\begin{env}[21.15.3]
\label{IV.21.15.3}
La structure de faisceau de groupes ordonnés sur $\mathscr{D}iv_X$ (21.1.6) induit
sur le sous-faisceau $\mathscr{D}iv_{X/S}$ une structure de faisceau de groupes
ordonnés, pour laquelle le faisceau de monoïdes des germes de sections positives est
$\mathscr{D}iv_X^+\cap\mathscr{D}iv_{X/S}$, que l'on note
$\mathscr{D}iv_{X/S}^+$. On a
$\Gamma(X,\mathscr{D}iv_{X/S}^+)=\operatorname{Div}^+(X)\cap
\operatorname{Div}(X/S)$; on note ce sous-monoïde de $\operatorname{Div}(X/S)$ par
$\operatorname{Div}^+(X/S)$, et il est formé des éléments $\geq0$ pour une structure
de groupe ordonné sur $\operatorname{Div}(X/S)$. Il résulte de (21.1.5.1) que
$\mathscr{D}iv_{X/S}^+$ est l'image dans $\mathscr{D}iv_{X/S}$ du sous-faisceau de
monoïdes
\[
\label{IV.21.15.3.1}
  \mathcal O_{X/S}^{\mathrm{reg}}=\mathcal O_X\cap\mathscr M_{X/S}^*.
  \tag{21.15.3.1}
\]
\end{env}

Pour tout ouvert $U$ de $X$, $\Gamma(U,\mathcal O_{X/S}^{\mathrm{reg}})$ est
l'ensemble des sections $t$ de $\mathcal O_X$ au-dessus de $U$ telles que $t$ soit
régulière et que $1/t$ appartienne à $\Gamma(U,\mathscr M_{X/S})$, ce qui signifie,
avec les notations de (20.6.1), que $t\in\mathscr T_{X/S}(U)$, de sorte que le faisceau
$\mathcal O_{X/S}^{\mathrm{reg}}$ n'est autre que le faisceau noté $\mathscr T_{X/S}$
dans (20.6.1). On peut donc écrire
\[
\label{IV.21.15.3.2}
  \mathscr{D}iv_{X/S}^+=\mathcal O_{X/S}^{\mathrm{reg}}/\mathcal O_X^*
  \tag{21.15.3.2}
\]
à un isomorphisme canonique près.

\begin{env}[21.15.3.3]
\label{IV.21.15.3.3}
Soit $D\in\operatorname{Div}^+(X/S)$ et considérons le sous-préschéma fermé $Y(D)$
de $X$ défini par l'Idéal $\mathscr J_X(D)$ de $\mathcal O_X$ (21.6.5); en vertu de
ce qui précède, pour tout $x\in Y(D)$, il y a un voisinage ouvert $U$ de $x$ dans $X$
et une section $t\in\mathscr T_{X/S}(U)$ telle que
$\mathscr J_X(D)|U=\mathcal O_U\cdot t$; puisque $x\in Y(D)$, l'image $t_{f(x)}$ de
$t$ dans $\Gamma(U\cap X_{f(x)},\mathcal O_{X_{f(x)}})$ appartient à l'idéal maximal
de $\mathcal O_{X_{f(x)},x}$ et en outre, par définition, pour tout $s\in S$, l'image
$t_s$ de $t$ dans $\Gamma(U\cap X_s,\mathcal O_{X_s})$ est une section
\emph{régulière}. On déduit donc de (11.3.8) et (19.2.4) que l'immersion canonique
$Y(D)\to X$ est \emph{transversalement régulière relativement à $S$ et de
codimension $1$} au point $x$. La réciproque étant immédiate, on voit qu'on peut
\emph{identifier canoniquement} les diviseurs positifs sur $X$ relatifs à $S$ aux
sous-préschémas fermés $Y$ de $X$ tels que l'injection canonique $Y\to X$ soit une
immersion transversalement régulière relativement à $S$ et de codimension $1$. Nous
ferons d'ordinaire cette identification.
\end{env}

\begin{proposition}[21.15.4]
\label{IV.21.15.4}
Soient $D$ un diviseur sur $X$, $\mathscr J_X(D)$ l'Idéal fractionnaire inversible
correspondant (21.2.5). Pour que $D\in\operatorname{Div}(X/S)$, il faut et il suffit
que, pour tout $x\in X$, on ait
$(\mathscr J_X(D)\cap\mathscr M_{X/S}^*)_x\ne0$ (ou, ce qui revient au même, que
l'on ait $\mathscr J_X(D)\subset\mathscr M_{X/S}$ et
$\mathscr J_X(D)^{-1}\subset\mathscr M_{X/S}$).
\end{proposition}

\begin{proof}
En effet, dire que $D\in\operatorname{Div}(X/S)$ signifie que pour tout $x\in X$, il
existe un voisinage ouvert $U$ de $x$ et une section
$f\in\Gamma(U,\mathscr M_{X/S}^*)$ tels que $D|U=\operatorname{div}_U(f)$; comme
$\mathscr J_X(D)|U$ est l'Idéal fractionnaire inversible $\mathcal O_Uf$, la
proposition en résulte aussitôt.
\end{proof}

\oldpage[IV-4]{331}
\begin{proposition}[21.15.5]
\label{IV.21.15.5}
Soient $D$ un diviseur sur $X$, $\mathcal O_X(D)$ l'Idéal fractionnaire inversible et
$s_D$ la section méromorphe régulière de $\mathcal O_X(D)$ définis canoniquement par
$D$ (21.2.8 et 21.2.9). Pour que $D\in\operatorname{Div}(X/S)$, il faut et il suffit
que $s_D\in\Gamma(X,\mathscr M_{X/S}(\mathcal O_X(D)))$.
\end{proposition}

\begin{proof}
En effet, si $U$ est un ouvert de $X$ tel que $D|U=\operatorname{div}_U(f)$, où
$f\in\Gamma(U,\mathscr M_X^*)$, dire que
$s_D\in\Gamma(U,\mathscr M_{X/S}(\mathcal O_X(D)))$ signifie, en vertu des
définitions (20.6.2), que $f^{-1}\in\Gamma(U,\mathscr M_{X/S})$, d'où la
proposition.
\end{proof}

L'interprétation des diviseurs sur $X$ à l'aide des classes
$\operatorname{cl}(\mathscr L,s)$ (21.2.11) permet donc d'interpréter les éléments de
$\operatorname{Div}(X/S)$ comme les couples (à isomorphisme près) $(\mathscr L,s)$
tels que $\mathscr L$ soit un $\mathcal O_X$-Module inversible et que $s$ soit une
section méromorphe de $\mathscr L$ au-dessus de $X$, \emph{régulière relativement à
$S$} (20.6.5, (iii)).

\begin{proposition}[21.15.6]
\label{IV.21.15.6}
Soit $D$ un diviseur sur $X$ relatif à $S$, et supposons que pour tout $x\in X$ tel
que $\operatorname{prof}(\mathcal O_{X_{f(x)},x})=1$, on ait $D_x\geq0$
(resp. $D_x=0$). Alors on a $D\geq0$ (resp. $D=0$).
\end{proposition}

\begin{proof}
Comme dans (21.1.8), on peut se borner au cas où $D=\operatorname{div}(\varphi)$, où
$\varphi$ est une fonction méromorphe régulière relative à $S$, et tout revient à voir
que si $D_x\geq0$ en tout point $x\in X$ tel que
$\operatorname{prof}(\mathcal O_{X_{f(x)},x})=1$, $\varphi$ est partout définie dans
$X$. Mais cette hypothèse signifie que, si $T=X-\operatorname{dom}(\varphi)$, on a
$\operatorname{prof}(\mathcal O_{X_{f(x)},x})\geq2$ pour tout $x\in T$, et il suffit
pour conclure d'appliquer (20.6.6).
\end{proof}

\begin{env}[21.15.7]
\label{IV.21.15.7}
Soit $X'$ un second $S$-préschéma plat et localement de présentation finie sur $S$, et
soit $g:X'\to X$ un $S$-morphisme. Si le $S$-morphisme $g$ est \emph{plat}, on sait
(21.4.5) que l'image réciproque par $g$ de tout diviseur sur $X$ est définie; si de
plus $D\in\operatorname{Div}(X/S)$, il résulte de la définition (21.15.2) et de
(20.6.8) que l'on a alors $g^*(D)\in\operatorname{Div}(X'/S)$.
\end{env}

\begin{env}[21.15.8]
\label{IV.21.15.8}
Soit $X'$ un second $S$-préschéma plat et localement de présentation finie sur $S$, et
soit $g:X'\to X$ un $S$-morphisme \emph{fini et plat}. Notons que $g$ est alors
nécessairement de présentation finie (1.4.3, 1.4.6 et 1.6.3), donc
$g_*(\mathcal O_{X'})$ est un $\mathcal O_X$-Module plat et de présentation finie, et
par suite \emph{localement libre} (2.1.12), autrement dit $g$ est un morphisme
localement libre (18.2.7); pour tout $s\in S$, le morphisme correspondant
$g_s:X'_s\to X_s$ est donc aussi fini et localement libre. On déduit alors de (21.5.2)
et (20.6.1) que pour tout ouvert $U$ de $X$ et toute section
$t'\in\Gamma(g^{-1}(U),\mathscr T_{X'/S})$, la norme $N_{X'/X}(t')$ appartient à
$\Gamma(U,\mathscr T_{X/S})$; le raisonnement de (21.5.3) prouve ensuite que pour tout
$\mathcal O_{X'}$-Module inversible $\mathscr L'$ et toute section méromorphe $u'$ de
$\mathscr L'$ au-dessus de $X'$, \emph{régulière relativement à $S$}, la norme
$N_{X'/X}(u')$ est une section méromorphe de
$\mathscr L=N_{X'/X}(\mathscr L')$ au-dessus de $X$, régulière relativement à $S$.
L'interprétation des diviseurs relatifs à $S$ donnée dans (21.15.5) et la définition de
l'image directe d'un diviseur (21.5.5) prouvent alors que \emph{pour tout diviseur}
$D'\in\operatorname{Div}(X'/S)$, on a $g_*(D')\in\operatorname{Div}(X/S)$.
\end{env}

\begin{env}[21.15.9]
\label{IV.21.15.9}
Considérons enfin un morphisme quelconque $S'\to S$, et (sous les hypothèses de
(21.15.1)) posons $X'=X\times_S S'$, qui est plat et localement de présentation finie
sur $S'$; si $p:X'\to X$ est la projection canonique, on a vu (20.6.9) qu'on a un
homomorphisme canonique
$p^*(\mathscr M_{X/S})\to\mathscr M_{X'/S'}$, qui transforme évidemment toute
section de $\mathscr M_{X/S}$ au-dessus d'un ouvert $U$, \emph{régulière relativement
à $S$}, en une section de $\mathscr M_{X'/S'}$ au-dessus de $p^{-1}(U)$, régulière
relativement à $S$ (20.6.5, (iii)); on conclut alors de la
\oldpage[IV-4]{332}
définition (21.15.2) et de l'exactitude à droite du foncteur $p^*$, que
l'homomorphisme précédent définit par passage aux quotients un homomorphisme canonique
\[
\label{IV.21.15.9.1}
  \operatorname{Div}(X/S)\longrightarrow\operatorname{Div}(X'/S')
  \tag{21.15.9.1}
\]
qui transforme évidemment les éléments de $\operatorname{Div}^+(X/S)$ en éléments de
$\operatorname{Div}^+(X'/S')$. On pose encore
$\operatorname{Div}(X'/S')=\operatorname{Div}_{X/S}(S')$ (resp.
$\operatorname{Div}^+(X'/S')=\operatorname{Div}_{X/S}^+(S')$), et on voit aussitôt
qu'on a ainsi défini deux foncteurs contravariants
\[
  \operatorname{Div}_{X/S}:\mathbf{Sch}_{/S}\longrightarrow\mathbf{Ab},\qquad
  \operatorname{Div}_{X/S}^+:\mathbf{Sch}_{/S}\longrightarrow\mathbf{Ens}
\]
de la catégorie des $S$-préschémas dans celle des groupes commutatifs (resp. des
ensembles). On verra plus loin (chap.~VI) des cas importants où le foncteur
$\operatorname{Div}_{X/S}^+$ est \emph{représentable} $(\mathbf 0_{\mathrm{III}},
8.1.8)$.

Pour tout diviseur $D\in\operatorname{Div}(X/S)$, l'image de $D$ par
l'homomorphisme (21.15.9.1) n'est autre d'ailleurs que l'\emph{image réciproque}
$p^*(D)$ (au sens de (21.4.2)) : l'existence de cette image réciproque et l'assertion
précédente sont en effet des conséquences immédiates de (20.6.5, (iii)) et (20.6.9).

Les éléments de $\operatorname{Div}(X'/S')$ s'appellent souvent, par abus de langage,
« \emph{familles de diviseurs sur $X$ relatifs à $S$, paramétrées par le
$S$-préschéma $S'$} »; on utilise surtout cette terminologie lorsqu'il s'agit de
diviseurs \emph{positifs}.
\end{env}
